\documentclass{article}
\usepackage{graphicx}
\usepackage[utf8]{inputenc}
\usepackage{array}
\usepackage[croatian]{babel} 
\usepackage{tabularx}
\usepackage{geometry}
\geometry{margin=2.5cm}
\usepackage{color}
\usepackage{titlesec}
\usepackage[table]{xcolor}
\usepackage{longtable}
\usepackage{xltabular}
\usepackage{enumitem}
\usepackage[hidelinks]{hyperref}
\usepackage{amsmath}
\setcounter{secnumdepth}{3}
\usepackage{array}
\usepackage{float}
\usepackage{booktabs}
\usepackage{colortbl}


\definecolor{headerblue}{RGB}{41, 128, 185}
\definecolor{lightgray}{RGB}{245, 245, 245}

\newcolumntype{L}{>{\raggedright\arraybackslash}X}
\newcolumntype{C}{>{\centering\arraybackslash}X}

\title{Informacioni sistem za pametnu farmu krava}
\author{Vedad Gaštan (Timski rad)}
\date{Oktobar 2025}

\begin{document}

\begin{titlepage}
    \centering
    \includegraphics[width=0.5\textwidth]{images/logo.png}\\[1em] 
    \vspace{3cm}

    \Huge\textbf{Informacioni sistem za pametnu farmu krava}\\[0.5em]
    \Huge\textbf{„OIS - Smart Cow Farm“}\par
    \vspace{3cm}

    \Large Projekat radili:\par
    \vspace{1em}
    \normalsize
    \begin{center}
    \begin{tabular}{c}
        Aldin Velić (vođa) \\
        Tarik Mujkić \\
        Danijal Alibegović \\
        Bakir Činjarević \\
        Emin Begić \\
        Enis Adilović \\
        Mirnes Fehrić \\
        Naila Delalić \\
        Vedad Gaštan \\
        Zana Beljuri
    \end{tabular}
    \end{center}
    \vspace{3cm}

    \normalsize Sarajevo, Oktobar 2025
\end{titlepage}

\tableofcontents

\newpage

\section{UVOD} % Section 1: UVOD (as in original)

\subsection{Svrha}
Svrha ovog projekta je razviti AI baziran informacioni sistem za pametnu farmu krava koji omogućava neinvazivno praćenje zdravstvenog stanja, produktivnosti i ponašanja svake krave, te automatizirano prikupljanje i obradu podataka radi unapređenja upravljanja farmom i optimizacije proizvodnje mlijeka.


\subsection{Konvencije dokumenta}
U dokumentu ćemo koristiti sljedeće akronime:

\begin{itemize}
 \item \textbf{IS} – Informacioni sistem
 \item \textbf{DB} – Baza podataka (Database)
 \item \textbf{GUI} – Grafički korisnički interfejs (Graphical User Interface)
 \item \textbf{SW} – Softver (Software)
 \item \textbf{AI} – Umjetna inteligencija (Artificial Intelligence)
 \item \textbf{IoT} – Internet stvari (Internet of Things) 
 \item \textbf{KM} – Konvertibilna marka
\end{itemize}

\subsection{Predviđeni korisnici i sugestije pri korištenju}

Predviđeni korisnici informacionog sistema pametne farme i njihove mogućnosti su:

\begin{itemize}
 \item \textbf{Osoblje farme} – koristi sistem za praćenje zdravstvenog stanja krava, evidenciju količine proizvedenog mlijeka, nadzor nad teladima i pravovremenu reakciju na upozorenja sistema. Obavezni su redovno ažurirati podatke o stanju životinja i uvjetima u štali kako bi se omogućila tačna analiza i podrška odlučivanju.

\item \textbf{Administrator sistema} – odgovoran je za održavanje, nadogradnju i sigurnost sistema. Ima pristup svim podacima, mogućnost upravljanja korisničkim nalozima, izmjene konfiguracija senzora i kontrolu nad cijelim infrastrukturnim okruženjem.

 \item \textbf{Partneri projekta} – visokoškolske i istraživačke institucije, veterinarske službe te kompanije uključene u razvoj sistema. Koriste sistem radi evaluacije funkcionalnosti, sprovođenja istraživanja i poboljšanja kvalitete upravljanja farmom.

 \item \textbf{Krajnji korisnici} – poljoprivrednici i menadžeri farmi koji koriste analitičke podatke sistema radi donošenja pametnih odluka, smanjenja troškova i povećanja produktivnosti. Dugoročno, dobit ostvaruju i krajnji potrošači kroz kvalitetnije i sigurnije proizvode.
\end{itemize}

\subsection{Opseg projekta}

Opseg ovog projekta obuhvata razvoj AI baziranog informacionog sistema za pametnu farmu krava koji omogućava automatizovano prikupljanje, obradu i praćenje ključnih podataka o stoci i uslovima na farmi. Cilj je olakšati svakodnevne procese upravljanja farmom kroz digitalizaciju i napredne tehnologije, uz povećanje efikasnosti proizvodnje mlijeka i unapređenje dobrobiti životinja.

Sistem će omogućiti:

\begin{itemize}
 \item Neinvazivno praćenje krava putem kamera i AI algoritama za prepoznavanje jedinki
 \item Evidentiranje proizvodnje mlijeka za svaku kravu i generisanje statistike produktivnosti
 \item Praćenje zdravlja i ponašanja radi pravovremenog otkrivanja bolesti ili stresa
 \item Nadzor uslova u štali putem IoT senzora (temperatura, vlaga, kvalitet zraka)
 \item Centralizovanu obradu podataka i prikaz u okviru web aplikacije dostupne osoblju i menadžmentu farme
\end{itemize}

Sistem će imati intuitivan korisnički interfejs i biti zasnovan na modernoj bazi podataka koja osigurava pouzdano skladištenje podataka. Platforma će podržavati više korisničkih uloga i omogućiti njihovu sigurnu autentifikaciju i autorizaciju.

Primarni fokus projekta je pružiti savremeno digitalno rješenje koje unapređuje produktivnost, smanjuje troškove poslovanja farme te omogućava održivu i efikasnu poljoprivrednu proizvodnju.

\subsection{Reference}

\begin{enumerate}
\item Tangorra, F. M., et al., “Internet of Things (IoT): Sensors Application in Dairy Cattle sector”, 2024. \textit{PMC}
 \item Posam, S., et al., “Automated Dairy Farm Management System Powered by IoT \& ML”, 2025. \textit{ScienceDirect}
 \item Palma, O., et al., “AI and Data Analytics in the Dairy Farms: A Scoping Review”, 2025. \textit{Animals, MDPI}
 \item Jothilakshmi, M. \& Illayabharathi, D., “Review on Tech Enabled Precision Dairy Farming – Case Study from Tamil Nadu, India”, 2024. \textit{International Journal of Veterinary Sciences and Animal Husbandry}
 \item Akbar, M.O., et al., “IoT for Development of Smart Dairy Farming”, 2020. \textit{Journal of Food Quality, Wiley Online Library}
 \item “Securing Smart Dairy Farms: A Cybersecurity Analysis of IoT”, 2024. \textit{AIP Conference Proceedings}
 \item “Dairy farming in the era of artificial intelligence: trend or a real game-changer”, 2024. \textit{Journal of Dairy Research}
 \item Zvanična web stranica Farme Spreča: \href{https://farmaspreca.ba/}{www.farmaspreca.ba}
 \item Inmedia članak \href{https://www.inmedia.ba/farma-spreca-najveca-farma-muznih-krava-u-bih-svakog-mjeseca-pola-miliona-litara-mlijeka/}{"Farma Spreča: Najveća farma muznih krava u BiH, svakog mjeseca pola miliona litara mlijeka"}
\end{enumerate}

\newpage
\section{HISTORIJA DOKUMENTA} % New section 2
\begin{center}
\begin{tabularx}{\textwidth}{|>{\centering\arraybackslash}p{3.5cm}|>{\centering\arraybackslash}p{2.5cm}|X|>{\centering\arraybackslash}p{1.5cm}|}
\hline
\rowcolor[gray]{0.9}
\textbf{Osoba} & \textbf{Datum} & \textbf{Razlog izmjene} & \textbf{Verzija} \\

\hline
Aldin Velić & 25.10.2025. & Inicijalna verzija dokumenta i uvod & 1.0 \\
\hline
Vedad Gaštan, Naila Delalić, Tarik Mujkić & 25.10.2025. & Analiza izvodivosti & 1.1 \\
\hline
Bakir Činjarević& 25.10.2025. & Upoznavanje s organizacijom & 1.2 \\
\hline
Danijal Alibegović& 25.10.2025. & Stakeholderi & 1.3 \\
\hline
Zana Beljuri, Mirnes Fehrić & 25.10.2025. & Misija, vizija, ciljevi organizacije & 1.4 \\
\hline
Enis Adilović & 25.10.2025. & Ciljevi projekta, Kvantitativni ciljevi, Kvalitativni ciljevi & 1.5 \\
\hline
Emin Begić & 25.10.2025. & Zahtjev za sistemom & 1.6 \\
\hline
Aldin Velić & 12.11.2025. & Definicija zahtjeva (FZ + NFZ)
 & 2.1 \\
\hline
Enis Adilović & 12.11.2025. & Slučajevi upotrebe (tabele)
 & 2.2 \\
\hline
Vedad Gaštan, Naila Delalić, Zana Beljuri, Mirnes Fehrić & 13.11.2025. & Dijagram slučajeva upotrebe, prepravka tabela slučajeva upotrebe
 & 2.3 \\
\hline
Bakir Činjarević & 13.11.2025. & Upitnik
 & 2.4 \\
\hline
Emin Begić & 13.11.2025. & Analiza dokumenata
 & 2.5 \\
\hline
Danijal Alibegović, Tarik Mujkić & 14.11.2025. & Intervju
 & 2.6 \\
\hline
Aldin Velić, Enis Adilović, Bakir Činjarević, Danijal Alibegović & 19.11.2025. & ERD dijagram
 & 3.1 \\
\hline
Vedad Gaštan & 22.11.2025. & Activity dijagram AD-1
 & 3.2 \\
\hline
Naila Delalić & 22.11.2025. & Activity dijagram AD-2
 & 3.3 \\
\hline
Tarik Mujkić & 23.11.2025. & Activity dijagram AD-3
 & 3.4 \\
 \hline
Emin Begić & 23.11.2025. & Activity dijagram AD-4
 & 3.5 \\
 \hline
Zana Beljuri & 23.11.2025. & Activity dijagram AD-5
 & 3.6 \\
\hline
Mirnes Fehrić & 23.11.2025. & Activity dijagram AD-6
 & 3.7 \\
 \hline

\hline
\end{tabularx}
\end{center}


\newpage
    \section{UPOZNAVANJE SA ORGANIZACIJOM} % New section 3
% This section is left empty as per request, to be filled in later.
\subsection{Opće informacije o organizaciji}
\textbf{Naziv organizacije:} Farma Spreča d.o.o.\\
\textbf{Adresa:} Donje Vukovije bb, 75260 Kalesija, Bosna i Hercegovina\\
\textbf{Djelatnost:} Proizvodnja mlijeka, uzgoj muznih krava, ratarstvo, proizvodnja bioplina\\
\textbf{Pravna forma:} Društvo sa ograničenom odgovornošću\\

\subsection{Kratka historija Farme Spreča}
Farma Spreča je najveća farma muznih krava u Bosni i Hercegovini, smještena na području općina Kalesije i Živinica, nadomak Tuzle. Farma je na ovom području postojala i prije rata, a nakon procesa privatizacije, farmu je otkupio Milkos i nakon ulaganja od gotovo 20 miliona KM, postala je glavni nosilac domaće proizvodnje mlijeka u BiH.

Farma Spreča posjeduje stado od oko 1600 krava uključujući sve kategorije, sa visokokvalitetnim grlima najboljih zdravstvenih i proizvođačkih karakteristika. Na farmi su visokokvalitetne mliječne krave pasmine Holstein uvezene iz najsavremenijih repro centara u Danskoj i Njemačkoj.

Farma danas posjeduje oko 1.500 stočnih grla, od čega su oko 800 muznih krava, koje mjesečno daju oko pola miliona litara mlijeka. Još 2012. godine farma je imala posjetu delegacije Evropske komisije, a Farma Spreča je jedna od prvih farmi u BiH koja je zadovoljila sve uslove za izvoz mlijeka u Evropsku uniju.

\subsection{Glavni proizvodi i usluge}
Primarna djelatnost - Proizvodnja mlijeka:
Farma Spreča ima kapacitet od 2500 grla, od čega 1500 muznih krava i 1000 različitih pratećih kategorija. Godišnja proizvodnja mlijeka iznosiće oko 12 miliona litara. Mlijeko u potpunosti otkupljuje sarajevska mljekara Milkos. Farma proizvodi visokokvalitetno mlijeko koje zadovoljava najstrožije kriterije zemalja EU zahvaljujući visokim higijenskim uslovima, pažljivoj ishrani i udobnom smještaju stoke.

Dodatne djelatnosti:
\begin{itemize}
    \item Ratarstvo: Farma Spreča obrađuje 1.200 hektara poljoprivrednog zemljišta, koje u većini služi za proizvodnju hrane za krave. Farma posjeduje vlastitu modernu mehanizaciju koja uključuje najsavremenije traktore, kombajne, poljoprivredne priključne mašine, prikolice, cisterne za stajnjak i mašine za spremanje stočne hrane.
    \item Reprodukcija: Krave se osjemenjavaju sjemenskim materijalom visokokvalitetnih bikova uz strogi nadzor i kontrolu stručnog veterinarskog kadra zaposlenog na farmi. 
    \item Planirani razvoj: Farma Spreča planira investicije u nove projekte u oblasti voćarstva i povrtlarstva i prerade voća i povrća, projekat obnovljivih izvora energije (biogasno postrojenje), izgradnju plastenika i staklenika, pogona za proizvodnju humusa, kao i najveći projekat u BiH - izgradnju moderne farme mliječnih koza kapaciteta 4.500 grla. 
\end{itemize}

\subsection{Tehnološka opremljenost}

Ukupan kapacitet sistema za mužu je oko 1.100 grla, što ostavlja prostor za dalji rast. Farma koristi savremenu opremu i tehnologiju:

\begin{itemize}
    \item Automatizovani sistemi muže
    \item Moderne štale sa kontrolisanim uslovima
    \item Savremena mehanizacija za obradu zemljišta
    \item Sistemi za upravljanje ishranom stoke
\end{itemize}

\subsection{Izazovi i potrebe za informatizacijom}

Uprkos modernoj opremi, Farma Spreča trenutno se suočava sa izazovima u:

\begin{itemize}
    \item Praćenju zdravstvenog stanja pojedinačnih grla u realnom vremenu
    \item Optimizaciji ishrane zasnovane na individualnim potrebama životinja
    \item Predviđanju proizvodnih kapaciteta i planiranju resursa
    \item Evidenciji reproduktivnih ciklusa i genetskih podataka
    \item Analizi podataka za donošenje poslovnih odluka
\end{itemize}

Upravo ove potrebe čine Farmu Spreča idealnim kandidatom za implementaciju pametnog sistema za upravljanje farmom zasnovanog na AI tehnologiji, koji će omogućiti:

\begin{itemize}
    \item Kontinuirani monitoring zdravlja krava
    \item AI-baziranu analitiku proizvodnih pokazatelja
    \item Automatsko praćenje i optimizaciju ishrane
    \item Prediktivno održavanje opreme
    \item Integrisanu evidenciju svih aspekata poslovanja farme
\end{itemize}

\newpage
\section{MISIJA, VIZIJA, CILJEVI ORGANIZACIJE} % New section 4

\subsection{Misija}
Misija naše organizacije je unaprijediti stočarsku proizvodnju kroz primjenu savremenih informacionih tehnologija i vještačke inteligencije.
Naš cilj je omogućiti farmerima jednostavan, precizan i neinvazivan način praćenja svake krave, proizvodnje mlijeka, zdravstvenog stanja i uslova u štali.
Digitalizacija procesa putem AI i IoT tehnologija omogućava donošenje odluka zasnovanih na podacima, čime se povećava produktivnost, smanjuju troškovi i unapređuje dobrobit životinja.

\subsection{Vizija}
Vizija naše organizacije je postati vodeći inovator u oblasti digitalne poljoprivrede u Bosni i Hercegovini i regiji.
Težimo stvaranju održivih, tehnološki naprednih farmi koje povezuju ljude, životinje i tehnologiju u jedinstven ekosistem.
Dugoročni cilj nam je da svaki farmer, bez obzira na veličinu proizvodnje, može upravljati svojom farmom putem inteligentnog sistema u realnom vremenu, povećavajući efikasnost i osiguravajući zdravlje životinja.

\subsection{Ciljevi}

Ciljevi organizacije podijeljeni su prema vremenskom okviru na kratkoročne, srednjoročne i dugoročne. \\

\textbf{Kratkoročni ciljevi (do 1 godine):}
\begin{itemize}
    \item Implementirati sistem za prepoznavanje i praćenje krava pomoću kamera i AI modela;
    \item Uspostaviti centralnu bazu podataka o proizvodnji mlijeka, zdravlju i ponašanju životinja;
    \item Omogućiti pristup web aplikaciji sa prikazom podataka, statistika i upozorenja u realnom vremenu;
    \item Povećati efikasnost proizvodnje mlijeka za najmanje 15\% u prvoj godini;
    \item Smanjiti stopu obolijevanja stoke i hitnih veterinarskih intervencija za 10\%.
\end{itemize}

\textbf{Srednjoročni ciljevi (1–5 godina):}
\begin{itemize}
    \item Integrisati sistem sa senzorima za kontrolu temperature, vlage i kvaliteta zraka u štali;
    \item Smanjiti troškove veterinarskih intervencija za 25\% kroz rano otkrivanje bolesti;
    \item Optimizirati ishranu krava analitikom podataka i smanjiti rasipanje hrane za 15\%;
    \item Unaprijediti reproduktivnu efikasnost stada za 10\%;
    \item Omogućiti integraciju sa nacionalnim bazama stočnog fonda i poljoprivrednim platformama.
\end{itemize}

\textbf{Dugoročni ciljevi (5+ godina):}
\begin{itemize}
    \item Proširiti sistem na druge vrste farmi (ovce, koze, perad);
    \item Implementirati pametni sistem na sve farme u okviru organizacije;
    \item Smanjiti ekološki otisak (potrošnja energije, vode i emisije) za 20\%;
    \item Postati regionalni centar za razvoj i izvoz rješenja u oblasti pametne poljoprivrede;
    \item Postati sinonim za održivu, digitalnu i etičku stočarsku proizvodnju.
\end{itemize}

\newpage
\section{STAKEHOLDERI}

\begin{itemize}
 \item \textbf{Administrator} – osoba koja održava sistem, briše, ažurira i dodaje dijelove sistema, te ima mogućnost uvida u sve parametre farme, uključujući zdravstveno stanje, proizvodnju mlijeka i podatke sa senzora. Pored održavanja sistema, ima i ulogu u sigurnosti podataka, posebno s obzirom na osjetljive zdravstvene informacije o kravama i farmama.

 \item \textbf{Osoblje farme} – koristi sistem za praćenje stanja krava, evidenciju proizvodnje mlijeka, zdravstvenog nadzora i brige o teladima. Također su odgovorni za ažuriranje stvarnih podataka i reagovanje na upozorenja sistema.

 \item \textbf{Kupci} – korisnici koji kupuju mliječne proizvode i druge proizvode farme. Sistem omogućava uvid u porijeklo mlijeka i transparentnost proizvodnog procesa, s obzirom da bi njihovi interesi mogli biti transparentnost, visoka kvaliteta, i etičko porijeklo proizvoda. 

 \item \textbf{Nevladine organizacije i Grupe za zaštitu okoliša} – Interesuju se za utjecaj sistema na dobrobit životinja (neinvazivni monitoring), te ekološku održivost i smanjenje otpada na farmi.

\item \textbf{Partneri projekta} – organizacije i institucije koje sarađuju na zajedničkom razvoju i unapređenju sistema. Potencijalni partneri pametne farme:
 \begin{itemize}
     \item Poljoprivredni fakultet Univerziteta u Sarajevu,
     \item Veterinarski fakultet Univerziteta u Sarajevu,
     \item Ministarstvo poljoprivrede, vodoprivrede i šumarstva FBiH,
     \item Poljoprivredne zadruge i lokalne farme,
     \item Agencija za sigurnost hrane BiH,
     \item Kompanije za razvoj AI i IoT tehnologija (Lanaco, Authority Partners, Symphony).
\end{itemize}

 \item \textbf{Sponzori projekta} – institucije i organizacije koje finansijski podržavaju razvoj i implementaciju sistema. Potencijalni sponzori:
 \begin{itemize}
 \item Grad Sarajevo,
 \item Federalno ministarstvo razvoja, poduzetništva i obrta,
 \item KULT za razvoj mladih (grantovi),
 \item UNDP BiH (projekti održivog razvoja),
 \item Evropska unija – fondovi za digitalizaciju poljoprivrede.
 \end{itemize}

 \item \textbf{Krajnji korisnici (beneficijari)} – poljoprivrednici i menadžeri farmi koji koriste sistem za praćenje zdravlja i produktivnosti stoke, optimizaciju troškova i povećanje efikasnosti. Dugoročno, beneficijari su i krajnji potrošači koji dobijaju zdravije i kvalitetnije proizvode.

\item \textbf{Vođa projekta i projektni tim}:\begin{itemize}
 \item Aldin Velić (vođa)
 \item Tarik Mujkić
 \item Danijal Alibegović
 \item Bakir Činjarević
 \item Emin Begić
 \item Enis Adilović
\item Mirnes Fehrić
 \item Naila Delalić
 \item Vedad Gaštan
 \item Zana Beljuri
 \end{itemize}
\end{itemize}

\newpage
\section{ZAHTJEV ZA SISTEMOM} % Section 6: ZAHTJEV ZA SISTEMOM (Original content)

\begin{table}[h!]
\renewcommand{\arraystretch}{1.4}
\setlength{\tabcolsep}{10pt}
\begin{tabularx}{\textwidth}{|>{\bfseries}l|X|}
\hline
\rowcolor[gray]{0.9}
\textbf{Element} & \textbf{Primjer} \\
\hline
Poslovna potreba: &
\begin{itemize}
 \item Praćenje zdravlja i produktivnosti svake krave
 \item Smanjenje troškova veterinarskih intervencija
 \item Optimizacija proizvodnje mlijeka
 \item Pravovremeno otkrivanje bolesti i stresa
 \item Automatizacija evidencije i smanjenje ručnog unosa podataka
\end{itemize} \\
\hline
Poslovni zahtjevi: &
\begin{itemize}
\item Automatsko prepoznavanje svake krave pomoću AI i kamera
 \item Evidencija količine proizvedenog mlijeka po kravi
 \item Praćenje promjena u ponašanju životinja
 \item Praćenje temperature, vlage i kvaliteta zraka u štali
 \item Web aplikacija za pregled podataka i generisanje izvještaja
\end{itemize} \\
\hline
Poslovna vrijednost: &
\begin{itemize}
 \item Povećanje proizvodnje mlijeka
 \item Smanjenje troškova liječenja i gubitaka
 \item Ušteda vremena zaposlenih kroz automatizaciju
 \item Poboljšanje dobrobiti životinja i kvaliteta proizvoda
\end{itemize} \\
\hline
Dodatni zahtjevi: &
\begin{itemize}
 \item Rok za implementaciju sistema: 01.03.2026.
 \item Sigurnost i privatnost podataka o farmi i životinjama
 \item Pouzdan rad sistema u realnom vremenu i u offline režimu
\end{itemize} \\
\hline

\end{tabularx}
\end{table}

\newpage
\section{CILJEVI PROJEKTA}

Cilj ovog projekta je dizajnirati i razraditi detaljan plan za razvoj informacionog sistema (IS) za pametnu farmu krava, koji će omogućiti neinvazivno praćenje stanja i performansi stada, podržati donošenje preciznih odluka na osnovu podataka i unaprijediti efikasnost i dobrobit životinja. Sistem će koristiti kombinaciju video-analitike, tehnika za prepoznavanje jedinki po jedinstvenim rasporedima fleka, senzora okoliša i centralizovane web aplikacije za upravljanje i vizualizaciju podataka.

\subsection{Kvantitativni ciljevi}

Kvantitativni ciljevi predstavljaju mjerljive aspekte sistema čiji uspjeh može biti objektivno procijenjen kroz konkretne metrike i brojčane pokazatelje.

\begin{itemize}
\item \textbf{AI model za prepoznavanje jedinki}\\
Razviti funkcionalni AI model za prepoznavanje krava sa tačnošću većom od 95\% na osnovu jedinstvenog rasporeda fleka. Model treba biti robustan u različitim uslovima osvjetljenja i uglovima snimanja, sa vremenom obrade manjim od 2 sekunde po identifikaciji.

\item \textbf{Praćenje proizvodnje mlijeka}\\
Implementirati sistem praćenja proizvodnje mlijeka sa mjernom greškom manjom od 2\%. Automatski evidentirati količinu proizvedenog mlijeka po kravi uz mogućnost agregacije podataka na dnevnom, sedmičnom i mjesečnom nivou sa preciznošću od najmanje 98\%.

\item \textbf{Integracija IoT senzora}\\
Uspostaviti integraciju sa svim IoT senzorima za praćenje uslova u štali (temperatura, vlaga, amonijak, CO$_2$) u realnom vremenu sa frekvencijom očitavanja od najmanje jednom u 5 minuta. Osigurati pouzdanost prenosa podataka veću od 99\%.

\item \textbf{Vrijeme odziva sistema}\\
Postići prosječno vrijeme odziva web aplikacije manje od 2 sekunde za standardne operacije i manje od 5 sekundi za kompleksne analitičke upite. Osigurati da sistem može istovremeno podržati najmanje 50 aktivnih korisnika bez degradacije performansi.

\item \textbf{Detekcija anomalija u ponašanju}\\
Razviti sistem za detekciju promjena u ponašanju krava sa najmanje 90\% tačnošću u identifikaciji potencijalnih zdravstvenih problema. Smanjiti vrijeme potrebno za identifikaciju bolesti ili stresa sa prosječnih 2-3 dana (manuelno praćenje) na manje od 12 sati.

\item \textbf{Smanjenje operativnih gubitaka}\\
Kroz implementaciju sistema ciljati smanjenje gubitaka uzrokovanih kasnom detekcijom bolesti za najmanje 30\% i povećanje prosječne proizvodnje mlijeka po kravi za 10-15\% kroz optimizaciju ishrane i uslova držanja.
\end{itemize}

\subsection{Kvalitativni ciljevi}

Kvalitativni ciljevi opisuju karakteristike sistema koje doprinose općoj kvaliteti, upotrebljivosti i dugoročnoj održivosti rješenja.

\begin{itemize}
\item \textbf{Intuitivan korisnički interfejs}\\
Omogućiti intuitivan i jednostavan GUI za osoblje farme koji ne zahtijeva prethodnu tehničku obuku. Dizajn aplikacije treba biti usklađen sa savremenim standardima korisničkog iskustva, sa jasnom navigacijom i pristupom svim ključnim funkcijama u najviše 3 klika.

\item \textbf{Pouzdanost i stabilnost sistema}\\
Osigurati visoku pouzdanost i stabilnost sistema sa minimalnim radnim vremenom od 99\%. Implementirati mehanizme za automatski oporavak od grešaka i redundantnost kritičnih komponenti kako bi se spriječio gubitak podataka i prekidi u radu. 

\item \textbf{Neinvazivno prepoznavanje jedinki}\\
Razviti koncept i algoritamske korake za neinvazivnu identifikaciju svake krave koristeći metode računalnog vida i mašinskog učenja. Definisati zahtjeve za kvalitet slike, optimalan raspored kamera u štali i kriterije za robusnost identifikacije u različitim uslovima (prljave krave, različito osvjetljenje, preklapanje životinja).

\item \textbf{Sveobuhvatno praćenje zdravlja i reprodukcije}\\
Definisati skup indikatora ponašanja koji mogu ukazivati na bolest, povredu ili stres (promjene u kretanju, izolacija od stada, smanjenje aktivnosti hranjenja, neobični zvukovi). Evidentirati reproduktivne događaje (estrusi, telenja) i zdravlje teladi uz automatsko detektovanje odstupanja od normalnih vrijednosti i obrazaca.

\item \textbf{Monitoring okoliša štale}\\
Specificirati integraciju senzora za temperaturu, vlagu i kvalitet zraka te način njihovog prostornog mapiranja unutar farme. Definisati pragove i politike alarma za okolinske parametre koji direktno utiču na zdravlje, produktivnost i dobrobit životinja. Omogućiti vizualizaciju podataka kroz toplinske mape i trendove.

\item \textbf{Podrška donošenju odluka}\\
Poboljšati procese donošenja odluka na farmi kroz analitičke izvještaje, prediktivne analize i automatske preporuke. Sistem treba pružiti jasne, razumljive i operativno korisne informacije koje farmerima omogućavaju brzu reakciju i dugoročno planiranje.

\item \textbf{Centralizovana obrada i vizualizacija podataka}\\
Dizajnirati web aplikaciju koja prikazuje sve relevantne podatke u centralizovanom formatu: detaljni profil svake krave sa historijom, trenutna i historijska proizvodnja, zdravstvena obavještenja, mape senzora i agregirane statistike na nivou stada. Osigurati mogućnost filtriranja, pretraživanja i prilagođenog izvještavanja.

\item \textbf{Sistem upozorenja i notifikacija}\\
Definisati mehanizme za automatsko generisanje upozorenja (e-mail, SMS, push notifikacije) pri detekciji anomalija ili kritičnih događaja. Implementirati prioritizaciju upozorenja prema ozbiljnosti problema i omogućiti personalizaciju postavki notifikacija prema potrebama korisnika.

\item \textbf{Sigurnost i privatnost podataka}\\
Osigurati visoku sigurnost podataka i usklađenost sa GDPR i sličnim regulativama o privatnosti (gdje je primjenjivo). Implementirati višeslojnu autentifikaciju korisnika, enkripciju podataka u prijenosu i skladištenju, te detaljne logove pristupa osjetljivim informacijama.

\item \textbf{Skalabilnost i modularnost}\\
Planirati arhitekturu koja omogućava skaliranje sistema bez značajnog pada performansi. Dizajn treba biti modularan kako bi se omogućilo lako dodavanje novih funkcionalnosti, integracija sa vanjskim servisima i prilagođavanje različitim veličinama farmi (od malih porodičnih do velikih komercijalnih).

\item \textbf{Operativna integracija i održavanje}\\
Navesti preporuke za integraciju sistema u postojeće farme i radne tokove osoblja uz minimalno narušavanje trenutnih operacija. Predložiti strategiju održavanja softvera i hardvera, procedure za ažuriranje AI modela sa novim podacima i plan obuke korisnika.

\item \textbf{Upravljanje zadacima i rasporedom}\\
Omogućiti planiranje i praćenje svakodnevnih aktivnosti na farmi kroz modul za upravljanje zadacima. Korisnici trebaju moći kreirati zadatke vezane za određene krave ili grupe (veterinarski pregledi, vakcinacije, tretmani), dodjeljivati ih članovima tima i pratiti njihovo izvršenje uz evidenciju utrošenih resursa i vremena.

\item \textbf{Povezanost sa vanjskim servisima}\\
Omogućiti integraciju sistema sa vanjskim servisima kao što su veterinarske ordinacije, mljekare i dobavljači. Automatizovati razmjenu relevantnih podataka, generisanje izvještaja za poslovne partnere i upravljanje lancima nabavke (zalihe hrane, automatsko naručivanje).

\item \textbf{Evaluacija i validacija}\\
Definisati metrike za mjerenje uspjeha implementacije (tačnost identifikacije, preciznost detekcije bolesti, smanjenje gubitaka, povećanje produktivnosti) i detaljan plan testiranja. Predložiti način prikupljanja povratnih informacija od krajnjih korisnika i proces iterativnog poboljšavanja rješenja.

\item \textbf{Etika, dobrobit životinja i održivost}\\
Uključiti etičke smjernice za neinvazivno praćenje i odgovornu upotrebu podataka, uz naglasak na dobrobit životinja kao primarni cilj. Razmotriti uticaj sistema na okoliš (potrošnja energije, elektronski otpad) i predložiti mjere za održivu primjenu tehnologije u skladu sa principima organske i etičke poljoprivrede.
\end{itemize}

\vspace{1em}
\noindent
Ovi ciljevi predstavljaju osnovni okvir za daljnju razradu arhitekture, funkcionalnih specifikacija i plana implementacije informacionog sistema. Projekt će se razvijati u skladu s principima modularnosti, mjerljivosti i etičnosti, s krajnjim ciljem da bude tehnički izvodiv, ekonomski opravdan, ekološki održiv i stvarno koristan krajnjim korisnicima---farmerima koji teže modernizaciji proizvodnje uz očuvanje dobrobiti životinja.

\newpage
\section{ANALIZA IZVODIVOSTI}

\subsection{Tehnička izvodljivost}

Projekat pametne farme krava sa integracijom vještačke inteligencije (AI) tehnički je izvodljiv, iako uključuje nekoliko izazova. Sistem predviđa upotrebu kamera, senzora za vlagu i mjerača količine mlijeka koji se povezuju s AI modelima za analizu zdravlja, ponašanja i proizvodnih parametara stoke. \\

\noindent
\textbf{Bliskost sa poslovnom oblašću – srednji rizik:}  
Tim koji razvija rješenje nema prethodnog direktnog iskustva u stočarskoj industriji, ali je kroz istraživanje tržišta, konsultacije sa stručnjacima i postojeće studije stekao solidno razumijevanje potreba farmi. Analizirani su ključni pokazatelji kao što su: navike u ishrani, fizičko kretanje, promjene u izgledu životinje, te varijacije u proizvodnji mlijeka — svi indikatori koji se mogu pratiti kroz AI.\\

\noindent
\textbf{Bliskost sa tehnologijom – nizak rizik:}  
Razvojni tim ima iskustva u radu s AI modelima, prepoznavanjem slika putem neuronskih mreža (CNN), obradom video streamova u realnom vremenu (OpenCV, YOLOv5), kao i radom sa IoT uređajima (ESP32, Arduino, Raspberry Pi). Također, poznate su platforme za obradu i pohranu podataka kao što su Firebase, PostgreSQL i InfluxDB.\\

\noindent
\textbf{Kompatibilnost sa postojećom tehnologijom – nizak rizik:}  
Predviđeni hardver (kamere, senzori za temperaturu, vlagu, mjerači mlijeka) koristi standardne protokole (npr. MQTT, HTTP REST, RTSP) te se može lako integrisati sa softverom koji će biti razvijen. Na strani korisnika (farmerski menadžment), koristiće se postojeći uređaji (laptopi, tableti, telefoni) bez potrebe za dodatnim ulaganjima.\\

\noindent
\textbf{Skalabilnost – srednji rizik:}  
Sistem je planiran tako da se može proširivati na veći broj krava i dodatne farme, ali to zahtijeva adekvatno planiranje arhitekture baze podataka i servera. Ako se ne osigura dovoljna procesorska snaga ili propusnost mreže, obrada video-streamova i senzorskih podataka može postati spora. U slučaju rasta broja uređaja, biće potrebno migrirati na cloud infrastrukturu (npr. AWS, Azure) i optimizovati AI modele za distribuirano procesiranje. Dakle, iako je sistem tehnički skalabilan, postoji rizik da će rast obima podataka povećati troškove održavanja i zahtijevati dodatne resurse.\\

\noindent
\textbf{Održavanje i nadogradnja – nizak do srednji rizik:}  
Sistem se planira modularno, što olakšava buduće nadogradnje (npr. dodavanje novih senzora, ažuriranje AI modela). Međutim, postoje rizici u kompatibilnosti verzija i potrebi za stručnim osobljem koje može ažurirati modele i softver bez prekida rada sistema. Ako farma nema stalnu IT podršku, održavanje može zavisiti od vanjskih partnera, što povećava rizik u pogledu troškova i vremena reakcije. Preporučuje se dokumentovan proces ažuriranja i redovni backup podataka kako bi se rizik sveo na minimum.\\

\noindent
\textbf{Veličina projekta – srednji rizik:}  
Za realizaciju je potreban tim od minimalno šest članova, uključujući AI inženjera, embedded programera, web/backend developera i osobu zaduženu za bazu podataka. Procijenjeno trajanje projekta je 4 do 6 mjeseci. Testiranje na terenu (na jednoj farmi) predviđeno je u finalnoj fazi projekta. Neophodna je bliska saradnja sa stakeholderima tokom razvoja.\\


\subsection{Ekonomska izvodivost}

\subsubsection{Troškovi razvoja sistema}

\textbf{Potrebna oprema:}
\begin{itemize}
  \item 4K kamere otporne na vanjske uslove (IP66 zaštita) - 15 kom
  \item Senzori za mjerenje temperature, vlažnosti i količine mlijeka - 100 kom
  \item Server za lokalnu pohranu podataka i backup (NAS) - 2 kom
  \item WiFi/LTE mrežna infrastruktura za prenos podataka
  \item UPS sistem za osiguranje napajanja
\end{itemize}

\begin{table}[h]
\centering
\caption{Troškovi razvoja}
\begin{tabular}{|l|r|}
\hline
\textbf{Vrsta troška} & \textbf{Iznos (KM)} \\ \hline
Razvojni tim & 120.000,00 \\ \hline
Troškovi hardvera i instalacije & 91.500,00 \\ \hline
Potrebni softveri & 3.000,00 \\ \hline
Računarska oprema & 5.000,00 \\ \hline
\textbf{Ukupno} & \textbf{219.500,00} \\ \hline
\end{tabular}
\end{table}

\subsubsection{Troškovi održavanja sistema}

\textbf{Varijabilni operativni troškovi:}
\begin{itemize}
    \item Električna energija: 15.330 kWh × 0,18 KM = 2.760 KM godišnje (rast 5\%)
    \item Internet konekcija: LTE/4G + backup = 1.560 KM godišnje
\end{itemize}

\begin{table}[h]
\centering
\caption{Troškovi održavanja}
\begin{tabular}{|l|r|r|r|r|r|r|}
\hline
\textbf{Vrsta troška} & \textbf{God 1} & \textbf{God 2} & \textbf{God 3} & \textbf{God 4} & \textbf{God 5} & \textbf{Ukupno} \\ \hline
Održavanje aplikacije & 5.000 & 5.000 & 5.000 & 5.000 & 5.000 & 25.000 \\ \hline
Tim za održavanje & 25.000 & 28.000 & 30.000 & 32.000 & 35.000 & 150.000 \\ \hline
Marketing & 2.500 & 2.500 & 2.500 & 2.500 & 2.500 & 12.500 \\ \hline
Stručno osoblje & 15.000 & 16.000 & 17.500 & 18.500 & 20.000 & 87.000 \\ \hline
Električna energija & 2.760 & 2.898 & 3.043 & 3.195 & 3.355 & 15.251 \\ \hline
Internet konekcija & 1.560 & 1.560 & 1.560 & 1.560 & 1.560 & 7.800 \\ \hline
\textbf{Ukupno} & \textbf{51.820} & \textbf{55.958} & \textbf{59.603} & \textbf{62.755} & \textbf{67.415} & \textbf{297.551} \\ \hline
\end{tabular}
\end{table}

\subsubsection{Određivanje toka novca}
\begin{table}[h]
\centering
\caption{Analiza toka novca}
\begin{tabular}{|l|r|r|r|r|r|r|r|}
\hline
\textbf{Pokazatelj} & \textbf{God 0} & \textbf{God 1} & \textbf{God 2} & \textbf{God 3} & \textbf{God 4} & \textbf{God 5} & \textbf{Ukupno} \\ \hline
Ukupan trošak & 219.500 & 51.820 & 55.958 & 59.603 & 62.755 & 67.415 & 517.051 \\ \hline
Ukupna korist & - & 50.000 & 100.000 & 150.000 & 180.000 & 200.000 & 680.000 \\ \hline
Neto korist & -219.500 & -1.820 & 44.042 & 90.397 & 117.245 & 132,585 & 162,949 \\ \hline
Kumulativni tok & -219.500 & -221.320 & -177.278 & -86,881 & 30,364 & 162,949 & - \\ \hline
\end{tabular}
\end{table}

\noindent
\textbf{\textit{Procjena priliva i odliva novca}}\\

\noindent
Na osnovu analize postojećih sistema za upravljanje farmama, očekuje se povećanje profita od 5\% u prvoj godini. Sistem omogućava dodatni profit od 50.000 KM u prvoj godini kroz bolje praćenje zdravlja životinja i optimizaciju ishrane. Taj iznos se progresivno povećava na 100.000 KM, 150.000 KM, 180.000 KM i 200.000 KM u narednim godinama.

\begin{table}[h]
\centering
\caption{Procjena troškova po godinama}
\begin{tabular}{|l|r|r|r|r|r|r|r|}
\hline
\textbf{Troškovi} & \textbf{God 0} & \textbf{God 1} & \textbf{God 2} & \textbf{God 3} & \textbf{God 4} & \textbf{God 5} & \textbf{Ukupno} \\ \hline
Održavanje aplikacije & - & 5.000 & 5.000 & 5.000 & 5.000 & 5.000 & 25.000 \\ \hline
Tim za održavanje & - & 25.000 & 28.000 & 30.000 & 32.000 & 35.000 & 150.000 \\ \hline
Marketing & - & 2.500 & 2.500 & 2.500 & 2.500 & 2.500 & 12.500 \\ \hline
Stručno osoblje & - & 15.000 & 16.000 & 17.500 & 18.500 & 20.000 & 87.000 \\ \hline
Električna energija & - & 2.760 & 2.898 & 3.043 & 3.195 & 3.355 & 15.251 \\ \hline
Internet & - & 1.560 & 1.560 & 1.560 & 1.560 & 1.560 & 7.800 \\ \hline
Razvojni tim & 120.000 & - & - & - & - & - & 120.000 \\ \hline
Hardver i instalacija & 91.500 & - & - & - & - & - & 91.500 \\ \hline
Softver i oprema & 8.000 & - & - & - & - & - & 8.000 \\ \hline
\textbf{Ukupno} & \textbf{219.500} & \textbf{51.820} & \textbf{55.958} & \textbf{59.603} & \textbf{62.755} & \textbf{67.415} & \textbf{517.051} \\ \hline
\end{tabular}
\end{table}

\begin{table}[h]
\centering
\caption{Procjena koristi po godinama}
\begin{tabular}{|l|r|r|r|r|r|r|}
\hline
\textbf{Koristi} & \textbf{God 1} & \textbf{God 2} & \textbf{God 3} & \textbf{God 4} & \textbf{God 5} & \textbf{Ukupno} \\ \hline
Dodatni profit farme & 50.000 & 100.000 & 150.000 & 180.000 & 200.000 & 680.000 \\ \hline
\end{tabular}
\end{table}

\subsubsection{Ocjenjivanje ekonomske vrijednosti projekta}

\textbf{ROI (Return on Investment):}

$$\text{ROI} = \frac{\text{Ukupna korist} - \text{Ukupan trošak}}{\text{Ukupan trošak}} \times 100\%$$

$$\text{ROI} = \frac{680.000 - 517.051}{517.051} \times 100\% = 31.51\%$$

\noindent
\textbf{BEP (Break-Even Point):}

\noindent
Projekat dostiže tačku pokrića između treće i četvrte godine. Kumulativni tok nakon treće godine iznosi -86,881 KM.

$$\text{BEP} = 3 + \frac{86,881}{30364 - (-86881)} = 3 + 0.74 = 3,74 \text{ godina}$$

$$\text{BEP} \approx 3 \text{ godine i } 9 \text{ mjeseci}$$

\noindent
\textbf{NPV analiza:} \\ \\
Diskontna stopa: 8\% \\
Formula za PV:
$$PV = \frac{FV}{(1 + r)^n}$$

\noindent
gdje je: FV = buduća vrijednost, r = diskontna stopa (0,08), n = broj godina\\

\noindent
Detaljni izračun PV koristi:

$$PV_1 = \frac{50.000}{(1,08)^1} = \frac{50.000}{1,08} = 46.296 \text{ KM}$$

$$PV_2 = \frac{100.000}{(1,08)^2} = \frac{100.000}{1,1664} = 85.734 \text{ KM}$$

$$PV_3 = \frac{150.000}{(1,08)^3} = \frac{150.000}{1,259712} = 119.074 \text{ KM}$$

$$PV_4 = \frac{180.000}{(1,08)^4} = \frac{180.000}{1,360489} = 132.301 \text{ KM}$$

$$PV_5 = \frac{200.000}{(1,08)^5} = \frac{200.000}{1,469328} = 136.117 \text{ KM}$$

\noindent
Detaljni izračun PV troška:

$$PV_1 = \frac{51.820}{(1,08)^1} = \frac{51.820}{1,08} = 47.981 \text{ KM}$$

$$PV_2 = \frac{55.958}{(1,08)^2} = \frac{55.958}{1,1664} = 47.975 \text{ KM}$$

$$PV_3 = \frac{59.603}{(1,08)^3} = \frac{59.603}{1,259712} = 47.315 \text{ KM}$$

$$PV_4 = \frac{62.755}{(1,08)^4} = \frac{62.755}{1,360489} = 46.127 \text{ KM}$$

$$PV_5 = \frac{67.415}{(1,08)^5} = \frac{67.415}{1,469328} = 45.882 \text{ KM}$$

\begin{table}[h]
\centering
\caption{NPV analiza}
\begin{tabular}{|l|r|r|r|r|r|r|r|}
\hline
\textbf{Pokazatelj} & \textbf{God 0} & \textbf{God 1} & \textbf{God 2} & \textbf{God 3} & \textbf{God 4} & \textbf{God 5} & \textbf{Ukupno} \\ \hline
Ukupna korist & - & 50.000 & 100.000 & 150.000 & 180.000 & 200.000 & 680.000 \\ \hline
PV koristi & - & 46.296 & 85.734 & 119.074 & 132.301 & 136.117 & 519.522 \\ \hline
Ukupni trošak & 219.500 & 51.820 & 55.958 & 59.603 & 62.755 & 67.415 & 517.051 \\ \hline
PV troška & 219.500 & 47.981 & 47.975 & 47.315 & 46.127 & 45.882 & 454.780 \\ \hline
\end{tabular}
\end{table}

\noindent
NPV (Net Present Value):

$$\text{NPV} = \sum PV_{\text{koristi}} - \sum PV_{\text{troška}}$$

$$\text{NPV} = 519.522 - 454.780 = 64.742 \text{ KM}$$


\noindent
Pošto je NPV $>$ 0, projekat obećava da će biti isplativ.\\

\noindent
\textbf{Finansijski pokazatelji (5 godina):}
\begin{itemize}
    \item ROI = 31,51\%
    \item BEP = 3,74 godina (3 godine i 9 mjeseci)
    \item NPV = 64.742 KM
\end{itemize}

\noindent
Projekat ostvaruje pozitivan ROI od 31,51\% u petogodišnjem periodu i dostiže tačku povrata nakon 3 godine i 9 mjesec. 


\subsection{Organizaciona izvodljivost}

Na organizacionu izvodljivost projekta pametne farme krava utiču faktori kao što su spremnost menadžmenta, komunikacija sa stakeholderima, stručnost zaposlenih i stepen prilagođavanja novim tehnologijama. Rizici su sljedeći:\\

\noindent
\textbf{Poklapanje ciljeva sa poslovnim planom – nizak rizik:}  
Menadžment farme iskazao je interes i spremnost da podrži digitalizaciju poslovanja. Projekat je u potpunosti u skladu sa strategijom modernizacije poljoprivrede, povećanja efikasnosti i smanjenja troškova proizvodnje. Implementacija informacionog sistema direktno doprinosi ciljevima održivog razvoja i poboljšanja dobrobiti životinja.\\

\noindent
\textbf{Komunikacija sa stakeholderima – nizak do srednji rizik:}  
Stakeholderi (osoblje farme, tehnički tim, sponzori i partneri) su uključeni u proces planiranja i dizajna sistema. Organizovani su redovni sastanci i konsultacije radi razmjene informacija i osiguranja transparentnosti. Rizik postoji ukoliko dođe do kašnjenja u komunikaciji između tehničkog tima i korisnika, što bi moglo uticati na pravovremenu isporuku funkcionalnosti.\\

\noindent
\textbf{Sposobnost menadžmenta – srednji rizik:}  
Menadžment posjeduje osnovno razumijevanje tržišta i digitalnih tehnologija, ali je potrebno dodatno usavršavanje u pogledu korištenja AI sistema i interpretacije analitičkih izvještaja. Obuka osoblja i menadžmenta planirana je u završnoj fazi projekta. Potrebno je osigurati kontinuitet obuka i tehničku podršku nakon implementacije.

\subsection{Operativna izvodljivost}
\textbf{Spremnost korisnika – srednji rizik:}  
Osoblje farme i menadžment su pokazali interes za primjenu digitalnih tehnologija, ali njihovo iskustvo u radu sa AI i IoT sistemima je ograničeno. Potrebno je organizovati praktične obuke i radionice kako bi se osiguralo da svi korisnici razumiju funkcionalnosti sistema i steknu povjerenje u njegov rad. Bez adekvatne edukacije, postoji rizik da će sistem biti korišten djelimično ili nepravilno.\\

\noindent
\textbf{Percepcija koristi – nizak rizik:}  
Korisnici će imati direktne koristi od sistema: lakše praćenje produktivnosti, smanjenje gubitaka i jednostavnije donošenje odluka. Sistem omogućava automatska upozorenja u slučaju zdravstvenih anomalija kod krava ili neadekvatnih uslova u štali, što značajno povećava operativnu vrijednost i motiviše korisnike na aktivno korištenje.





\subsection{Zakonska izvodljivost}

Sistem za pametnu farmu krava nije u sukobu sa zakonima Bosne i Hercegovine, ali zahtijeva usklađenost sa relevantnim pravnim propisima. Ključni zakoni koji se moraju poštovati su:
\begin{itemize}
    \item \textbf{Zakon o zaštiti ličnih podataka BiH} (``Službeni glasnik BiH'', br. 49/06, 76/11 i 89/11),
    \item \textbf{Zakon o zaštiti i dobrobiti životinja} (``Službeni glasnik BiH'', br. 25/09),
    \item \textbf{Zakon o zaštiti na radu} (``Službene novine FBiH'', br. 79/20),
    \item te drugi propisi koji uređuju oblast video-nadzora i sigurnosti informacija.
\end{itemize}

\noindent
\textbf{Privatnost i video-nadzor – srednji rizik:}  
Sistem koristi video kamere za praćenje stanja životinja. Ukoliko kamere pokrivaju područja u kojima borave radnici, potrebno je jasno označiti da je prostor pod video-nadzorom i obavijestiti zaposlene o svrsi snimanja, u skladu sa Zakonom o zaštiti ličnih podataka. Snimci se smiju koristiti isključivo za praćenje dobrobiti životinja i sigurnost objekta, te se moraju čuvati ograničen vremenski period (npr. 30 dana), nakon čega se automatski brišu.\\

\noindent
\textbf{Zaštita podataka i sigurnost – nizak rizik:}  
Svi podaci prikupljeni putem senzora i kamera moraju biti zaštićeni enkripcijom i pohranjeni na siguran server. Pristup podacima mora biti ograničen samo na ovlaštene osobe, uz primjenu autentifikacije i kontrolu pristupa. Ukoliko se koristi cloud pohrana, potrebno je osigurati da se podaci nalaze na serverima koji ispunjavaju sigurnosne standarde EU (npr. GDPR-kompatibilni data centri).\\

\noindent
\textbf{Dobrobit životinja – nizak rizik:}  
Zakon o zaštiti i dobrobiti životinja nalaže da životinje moraju biti zaštićene od nepotrebne patnje, povreda i stresa. Predloženi sistem doprinosi poštivanju tog zakona, jer omogućava rano otkrivanje bolesti i nadzor uslova u štali (temperatura, vlaga, kvalitet zraka). Senzori i uređaji koji se koriste moraju biti certificirani i ne smiju ugrožavati fizičko ili psihičko stanje životinje.\\

\newpage
\section{INTERVJU}

\subsection*{Intervju 1 — Tehnička perspektiva}

\noindent
\textbf{Intervjuisao:} Danijal Alibegović \\
\textbf{Datum/Vrijeme:} 14.11.2025 / 21:28 \\
\textbf{Mjesto:} Farma Spreča / Sarajevo \\
\textbf{Tema:} Tehnički zahtjevi i arhitektura OIS -- Smart Cow Farm \\
\textbf{Ispitanik:} IT inžinjer (razvoj i implementacija)

\vspace{1em}
\noindent\textbf{Ukupno trajanje:} 33 minute

\vspace{0.5em}
\begin{xltabular}{\textwidth}{|c|L|L|}
\hline
\rowcolor{headerblue}
\textcolor{white}{\textbf{Vrijeme}} & 
\textcolor{white}{\textbf{Pitanje}} & 
\textcolor{white}{\textbf{Odgovor}} \\
\hline
\endfirsthead

\hline
\rowcolor{headerblue}
\textcolor{white}{\textbf{Vrijeme}} & 
\textcolor{white}{\textbf{Pitanje}} & 
\textcolor{white}{\textbf{Odgovor}} \\
\hline
\endhead

\hline
\endfoot

\rowcolor{lightgray}
3 min &
\textbf{Koje su ključne komponente sistema i njihova uloga?} &
Sistem se sastoji od tri glavna sloja. Video-analitika koristi kamere i AI modele za neinvazivnu identifikaciju i praćenje ponašanja krava. IoT senzori prate okolišne uslove (temperatura, vlaga, NH$_3$, CO$_2$) i podatke s mljekarskih aparata. Centralna baza podataka i web aplikacija omogućavaju sigurno čuvanje, vizualizaciju, historiju i upravljanje alarmima.

Komponente su povezane kroz standardne protokole (RTSP, MQTT, REST) i rade kao jedinstven nadzorni sistem.
 \\
\hline
3 min &
\textbf{Koje performanse očekujete od AI modela za prepoznavanje krava?} &
Model mora imati tačnost $\ge95\%$ na test setu, raditi stabilno u različitom osvjetljenju i iz različitih uglova, te izvršiti identifikaciju u $\le2$ sekunde. Osim toga, očekuje se otpornost na prljavštinu, djelimičnu blokadu kamere i promjene izgleda životinje tokom vremena. Ručni override ostaje dostupan radi sigurnosti.\\
\hline
\rowcolor{lightgray}
3 min &
\textbf{Kako integrišete IoT senzore i koja je frekvencija očitanja?} &
IoT uređaji se povezuju putem MQTT ili REST protokola prema centralnom brokeru/serveru. Senzori šalju očitanja svakih 5 minuta, uz dnevni gubitak paketa manji od 1\%. Svaki senzor ima definisane pragove, politiku alarma i automatsku ponovnu kalibraciju. Sistem čuva historiju promjena radi audit traila.\\
\hline
4 min &
\textbf{Kako će se pratiti zdravlje i ponašanje krava?} &
Koristimo kombinaciju video-analitike i IoT signala: modeli detektuju promjene u kretanju, držanju, apetitu i abnormalnim obrascima ponašanja. IoT senzori dopunjuju sliku praćenjem okoliša. Sistem treba prepoznati anomaliju u <12 sati za $\ge90\%$ slučajeva, uz <5\% lažnih alarma. Svaka detekcija automatski se bilježi u profil životinje. \\
\hline
\rowcolor{lightgray}
3 min &
\textbf{Kako funkcionišu upozorenja/notifikacije i prioriteti?} &
Sistem generiše push, e-mail ili SMS notifikacije u $\le60$ sekundi nakon događaja. Postoje tri nivoa prioriteta: informativno, visoko i hitno. Korisnici mogu birati kanal i prag osjetljivosti, a svi alarmi se čuvaju u logu sa timestampom i kontekstom (krava, senzor, zona). \\
\hline
4 min &
\textbf{Koje sigurnosne mjere primjenjujete u prijenosu i pohrani?} &
Koristimo TLS 1.2+ za sav mrežni promet, HSTS za zaštitu browser pristupa, AES-256 za enkripciju podataka u mirovanju i bcrypt/Argon2 za hashiranje lozinki. Za administratore je obavezan 2FA. Sistem uključuje audit logove, kontrolu pristupa po rolama i politiku jakih lozinki.\\
\hline
\rowcolor{lightgray}
3 min &
\textbf{Gledajući na backup i oporavak, koje RPO/RTO ciljeve ciljate?} &
Implementiramo dnevne full i inkrementalne backup-e s ciljem RPO $\le$ 15 min i RTO $\le$ 4h. Backup se automatski verifikuje i čuva na odvojenoj lokaciji, a procedura oporavka se testira kvartalno. \\
\hline
3 min &
\textbf{Kakva je skalabilnost i performanse UI-a?} &
UI i backend moraju podržati $\ge$50 simultanih korisnika i obradu podataka za 1500+ krava bez degradacije. Standardne operacije (pregled, filtriranje, dashboard) moraju biti brže od 2s, dok napredne analitike moraju završiti za manje od 5s. Sistem je predviđen za horizontalno skaliranje po potrebi.\\
\hline
\rowcolor{lightgray}
3 min &
\textbf{Kako radi offline režim?} &
Terenski uređaji keširaju podatke lokalno i izvršavaju rad bez mreže. Po ponovnom spajanju, automatski se pokreće sinhronizacija u roku $\le$15 min, uz rješavanje konflikata kroz pregled promjena. Offline režim uključuje osnovne funkcije kao unos i pregled zadataka.\\
\hline
4 min &
\textbf{Kako ćete mjeriti kvalitet AI i planirati re-obuku?} &
Kvalitet se mjeri kroz metrike: identifikacija $\ge95\%$, zdravlje/ponašanje $\ge90\%$, lažni alarmi <5\%. Periodično se skuplja novi dataset sa farme, a kompletna re-obuka modela radi se npr. svaka 3 mjeseca ili pri većim odstupanjima u performansama. Postoji definisana procedura A/B evaluacije novog modela. \\
\hline
\end{xltabular}

\subsection*{Intervju 2 — Operativna perspektiva}

\noindent
\textbf{Intervjuisao:} Tarik Mujkić \\
\textbf{Datum/Vrijeme:} 13.11.2025. / 13:12 \\
\textbf{Mjesto:} Farma Spreča / Sarajevo \\
\textbf{Tema:} Svakodnevna upotreba sistema i operativni tokovi \\
\textbf{Ispitanik:} Član osoblja farme (zdravstveni nadzor)

\vspace{1em}
\noindent\textbf{Ukupno trajanje:} 34 minute

\vspace{0.5em}
\begin{tabularx}{\textwidth}{|c|L|L|}
\hline
\rowcolor{headerblue}
\textcolor{white}{\textbf{Vrijeme}} & 
\textcolor{white}{\textbf{Pitanje}} & 
\textcolor{white}{\textbf{Odgovor}} \\
\hline
\rowcolor{lightgray}
3 min &
\textbf{Kako upisujete novu kravu i kako je sistem prepoznaje kasnije?} &
Unosimo osnovne podatke i nekoliko referentnih fotografija. Sistem generiše jedinstveni ID i QR oznaku, a kasnija identifikacija se vrši automatski preko kamere korištenjem prepoznavanja šara dlake. \\
\hline
3 min &
\textbf{Kako evidentirate proizvodnju mlijeka po muži?} &
Podaci se automatski povlače sa muzačkih aparata i senzora sa tačnošću mjerenja do $\leq 2\%$. Po potrebi je moguć ručni unos koji se posebno označi u evidenciji. \\
\hline
\rowcolor{lightgray}
4 min &
\textbf{Šta radite kada sistem detektuje pad proizvodnje veći od 15\%?} &
Sistem generiše hitan alarm, automatski označava kravu i navodi moguće uzroke. Obavještava se veterinar, bilježi se intervencija i prema potrebi koriguje plan ishrane. \\
\hline
3 min &
\textbf{Kako koristite trendove i prognoze u dashboardu?} &
Dashboard prikazuje dnevne, sedmične i mjesečne trendove te prognoze rasta ili pada proizvodnje. Podatke izvozimo u PDF ili CSV jednim klikom za potrebe analize ili slanja upravi. \\
\hline
\rowcolor{lightgray}
4 min &
\textbf{Kako postupate po zdravstvenim upozorenjima (npr. estrus, stres)?} &
Veterinar automatski dobija notifikaciju sa ID-om krave, opisom simptoma i kratkim videosnimkom. Može potvrditi ili odbaciti upozorenje, a sistem potom prati terapiju i dalji status životinje. \\
\hline
4 min &
\textbf{Kako koristite modul zadataka u svakodnevnoj organizaciji posla?} &
Zadaci se kreiraju i dodjeljuju radnicima, uz automatske notifikacije i praćenje vremena i utrošenih resursa. Rutinski poslovi poput hranjenja i čišćenja generišu se automatski prema definisanim rasporedima. \\
\hline
\rowcolor{lightgray}
3 min &
\textbf{Koje izvještaje dobijate i kome ih šaljete?} &
Sistem generiše standardne izvještaje o proizvodnji, zdravlju i reprodukciji za manje od 30 sekundi. Izvještaji mogu biti zakazani za automatsko slanje menadžmentu ili veterinarima. \\
\hline
3 min &
\textbf{Koliko je interfejs jednostavan i kakva vam je obuka potrebna?} &
Dizajniran je tako da 90\% svakodnevnih zadataka bude obavljeno u najviše 3 klika. Planirana je kratka praktična obuka na licu mjesta, uz interaktivne tutorijale u samoj aplikaciji. \\
\hline
\rowcolor{lightgray}
3 min &
\textbf{Šta radite kada nema interneta?} &
Aplikacija radi offline i omogućava unos svih podataka. Čim se veza vrati, sistem automatski sinhronizuje sve unose u roku od najviše 15 minuta. \\
\hline
4 min &
\textbf{Kako se poštuje privatnost zaposlenih i video-nadzor?} &
Svi nadzirani prostori su jasno označeni, a svrha video-nadzora transparentno komunicirana. Snimci se čuvaju ograničen vremenski period, na primjer 30 dana, uz strogu kontrolu pristupa. \\
\hline
\end{tabularx}


\newpage

% =======================================================
% INTERVJU 3 — UPRAVLJANJE I FINANSIJE (MENADŽER)
% =======================================================
\subsection*{Intervju 3 — Upravljanje i finansije}

\noindent
\textbf{Intervjuisao:} Danijal Alibegović \\
\textbf{Datum/Vrijeme:} 15.11.2025 / 16:06 \\
\textbf{Mjesto:} Uprava farme / Sarajevo \\
\textbf{Tema:} Budžet, koristi, ROI i plan digitalizacije \\
\textbf{Ispitanik:} Menadžer farme / finansijski direktor

\vspace{1em}
\noindent\textbf{Ukupno trajanje:} 34 minute

\vspace{0.5em}
\begin{tabularx}{\textwidth}{|c|L|L|}
\hline
\rowcolor{headerblue}
\textcolor{white}{\textbf{Vrijeme}} & 
\textcolor{white}{\textbf{Pitanje}} & 
\textcolor{white}{\textbf{Odgovor}} \\
\hline
\rowcolor{lightgray}
3 min &
\textbf{Koji je okvirni budžet za razvoj?} &
Ukupan početni trošak iznosi 219.500 KM: razvojni tim 120.000 KM, hardver i instalacija 91.500 KM, softver 3.000 KM i računarska oprema 5.000 KM. \\
\hline
4 min &
\textbf{Koji su glavni stavovi opreme?} &
Oprema uključuje: 15 IP66 4K kamera, oko 100 senzora (temperatura, vlaga, količina mlijeka), 2 NAS servera, mrežnu infrastrukturu i UPS. \\
\hline
\rowcolor{lightgray}
4 min &
\textbf{Koliki su očekivani godišnji troškovi održavanja?} &
Troškovi održavanja: 51.820 KM u prvoj godini, zatim rastu do 67.415 KM.  
Ukupan petogodišnji trošak: 297.551 KM.  \\
\hline
4 min &
\textbf{Kakve finansijske koristi očekujete i kada je BEP?} &
Ukupna korist za pet godina: 680.000 KM.  
ROI = 31,51\%.  
NPV = 64.742 KM (diskont 8\%).  
BEP nakon 3,74 godine. \\
\hline
\rowcolor{lightgray}
2 min &
\textbf{Rok implementacije i potreban tim?} &
Realizacija traje 4--6 mjeseci uz tim od 6+ članova.  \\
\hline
3 min &
\textbf{Koja je poslovna vrijednost digitalizacije?} &
Povećanje proizvodnje, ranija detekcija bolesti, manji gubici, optimizacija ishrane i bolja efikasnost rada, što donosi dodatni profit od 50k do 200k KM godišnje. \\
\hline
\rowcolor{lightgray}
3 min &
\textbf{Koje KPI-eve ćete pratiti?} &
Povećanje proizvodnje po kravi 10--15\%, smanjenje gubitaka $\ge30\%$, vrijeme reakcije na alarme, sistemski uptime $\ge99\%$.  \\
\hline
4 min &
\textbf{Glavni rizici i mjere ublažavanja?} &
Rizici uključuju trošak obrade video-streamova koji se rješava cloud skaliranjem, kontinuitet obuke korisnika, te sigurnost.
Mjere ublažavanje koje planiramo koristiti su: TLS, AES-256, 2FA, redovni backupi i monitoring opreme. \\
\hline
\rowcolor{lightgray}
3 min &
\textbf{Usklađenost sa zakonima i privatnošću?} &
Poštovanje zaštite ličnih podataka, dobrobiti životinja i video-nadzora; ograničena retencija snimaka (npr. 30 dana). \\
\hline
3 min &
\textbf{Plan širenja i integracija?} &
Otvoren REST API za integracije (mljekara ERP, veterinarski IS); modularna arhitektura omogućava budući rast i skaliranje. \\
\hline
\end{tabularx}

\newpage
\section{UPITNIK ZA ANALIZU ZAHTJEVA}

U nastavku slijedi primjer upitnika kojim bi se analizirala potreba za ovakvim informacionim sistemom i stavovima korisnika koji bi njime upravljali.
\\
\rule{\linewidth}{0.4pt}
\\

Poštovani,

U okviru razvoja informacionog sistema za pametnu farmu krava, molimo Vas da izdvojite nekoliko minuta za popunjavanje ovog upitnika. Vaše mišljenje i iskustvo su od ključnog značaja za razvoj sistema koji će najbolje odgovarati potrebama moderne farme. Upitnik je anoniman, a prikupljeni podaci će se koristiti isključivo u svrhu analize zahtjeva za razvoj sistema.

Zahvaljujemo Vam na saradnji!

\subsection{Opšti podaci}

\textbf{1. Vaša uloga na farmi:} Vlasnik farme, Upravitelj/menadžer, Veterinar, Radnik na farmi, Drugo
\\
\textbf{2. Veličina farme (broj grla):} Do 50 krava, 51-100 krava, 101-200 krava, 201-500 krava, Više od 500 krava
\\
\textbf{3. Koliko dugo se bavite stočarstvom?} Do 2 godine, 3-5 godina, 6-10 godina, 11-20 godina, Više od 20 godina

\subsection{Trenutno stanje i izazovi}

\textbf{4. Kako trenutno vodite evidenciju o kravama na farmi?} Ručno (papir, bilježnice), Excel/Google Sheets, Specijalizovani softver, Kombinacija navedenog, Ne vodim detaljnu evidenciju
\\
\textbf{5. Koji su trenutno najveći izazovi u upravljanju farmom? (možete odabrati više odgovora)} Praćenje zdravstvenog stanja životinja, Praćenje proizvodnje mlijeka, Organizacija ishrane, Praćenje reprodukcije, Evidencija veterinarskih intervencija, Administracija i dokumentacija, Nedostatak vremena za sve poslove, Drugo
\\
\textbf{6. Koliko često se susrećete sa problemom kasnog otkrivanja bolesti kod krava?} Vrlo često, Često, Povremeno, Rijetko, Nikada

\subsection{Potrebe i očekivanja od novog sistema}

\textbf{7. Koje podatke o kravama smatrate najvažnijim za praćenje? (rangirati od 1-5, gdje je 1 najvažnije)} Zdravstveno stanje, Proizvodnja mlijeka, Reprodukcija i telenja, Ishrana, Ponašanje i aktivnost
\\
\textbf{8. Koje funkcionalnosti bi Vam bile najkorisnije u novom sistemu? (možete odabrati više odgovora)} Automatsko praćenje zdravstvenog stanja krava, Rano upozorenje na potencijalne zdravstvene probleme, Evidencija proizvodnje mlijeka po kravi, Praćenje ciklusa teljenja i osemenjavanja, Digitalna istorija svake krave, Automatsko generisanje izvještaja, Praćenje troškova po kravi, Upravljanje zalihama hrane i lijekova, Mobilna aplikacija za pristup podacima, Drugo
\\
\textbf{9. Koliko bi Vam bilo važno da sistem koristi AI za analizu i preporuke?} Veoma važno, Važno, Neutralno, Nije prioritet, Nepotrebno
\\
\textbf{10. Koje parametre biste željeli da AI analizira kod krava? (možete odabrati više odgovora)} Temperatura tijela, Aktivnost i kretanje, Apetit i unos hrane, Proizvodnja mlijeka, Ponašanje, Fizički izgled, Drugo

\subsection{Tehnički aspekti}

\textbf{11. Koji uređaji su dostupni na Vašoj farmi? (možete odabrati više odgovora)} Računar, Tablet, Pametni telefon, Internet konekcija, Kamere, Senzori (temperatura, vlažnost, itd.), Nemam ništa od navedenog
\\
\textbf{12. Koliko često koristite tehnologiju?} Veoma često koristim tehnologiju, Umjereno često, Osnovno znanje, Rijetko koristim tehnologiju, Ne koristim tehnologiju
\\
\textbf{13. Koliko bi sistem trebao biti jednostavan za korištenje?} Maksimalno jednostavan, bez tehničkih detalja, Relativno jednostavan, sa osnovnim funkcijama, Može imati napredne funkcije, ako su dobro objašnjene, Preferiram kompleksne sisteme sa mnogo opcija

\subsection{Implementacija i podrška}

\textbf{14. Koliko brzo biste željeli da se naviknete na novi sistem?} Odmah, uz minimalnu obuku, U roku od nekoliko dana, U roku od nekoliko sedmica, Spreman/na sam uložiti više vremena za obuku
\\
\textbf{15. Koja vrsta podrške bi Vam bila potrebna? (možete odabrati više odgovora)} Video tutorijali, Pisana uputstva, Telefonska podrška, Podrška uživo/na licu mjesta, Online chat podrška, Redovne obuke
\\
\textbf{16. Koliko često biste željeli primati izvještaje o stanju farme?} Dnevno, Sedmično, Mjesečno, Po potrebi/na zahtjev, Ne trebaju mi automatski izvještaji

\subsection{Dodatna pitanja}

\textbf{17. Koje dodatne funkcionalnosti biste željeli vidjeti u sistemu?}
\\
\textbf{18. Postoje li specifični problemi na Vašoj farmi koje biste željeli da novi sistem riješi?}
\\
\textbf{19. Da li biste bili zainteresovani za učešće u testiranju sistema prije zvanične implementacije?} Da, Ne, Možda
\\
\textbf{20. Dodatni komentari ili sugestije}

\subsection{Kontakt informacije (opcionalno)}

Ako želite da Vas kontaktiramo sa dodatnim pitanjima ili rezultatima analize, molimo Vas da ostavite svoje kontakt podatke.

\vspace{0.5cm}

\textit{Hvala Vam što ste izdvojili vrijeme za popunjavanje upitnika!}

\vspace{1em}
\hspace{-1.75em}
\rule{\linewidth}{0.4pt}

\newpage

\section{ANALIZA DOKUMENATA}

U nastavku su opisani ključni dokumenti koji nastaju u radu pametne farme i koje predloženi informacioni sistem treba da podrži, evidentira ili automatski generiše.

\subsection*{Karton krave (evidencija krava)}
Ovaj dokument predstavlja centralni dosije svake krave u stadu i koristi se pri registraciji, praćenju proizvodnje i zdravlja. \\

\noindent Elementi dokumenta su sljedeći:
\begin{itemize}
    \item Jedinstveni ID krave (AI/QR kod, ušni broj)
    \item Ime / oznaka krave
    \item Rasa
    \item Datum rođenja i datum dolaska na farmu
    \item Porijeklo (rođena na farmi / kupljena, ID majke)
    \item Trenutni status (u laktaciji, steona, zasušena, telad, otpisana)
    \item Početna težina i procijenjena trenutna težina
    \item Kratak opis izgleda i referentne fotografije (šara fleka)
    \item Prosječna dnevna proizvodnja mlijeka
    \item Napomene (specifičnosti, alergije, posebni tretmani)
\end{itemize}

\subsection*{Karton zdravstvenog stanja krave}
Za svaku kravu vodi se detaljan karton zdravstvenog stanja koji objedinjuje podatke koje je detektovao AI sistem i intervencije veterinara. \\

\noindent Elementi dokumenta su sljedeći:
\begin{itemize}
    \item Jedinstveni ID krave
    \item Datum i vrijeme pregleda ili upozorenja sistema
    \item Razlog otvaranja slučaja (AI detekcija, vizuelni pregled, pad proizvodnje i sl.)
    \item Opis simptoma i poremećaja ponašanja
    \item Rezultati AI analize (tip anomalije, nivo rizika)
    \item Dijagnoza veterinara
    \item Propisana terapija (naziv lijeka, doza, trajanje)
    \item Evidencija primjene terapije (datumi, odgovorno osoblje)
    \item Plan kontrolnih pregleda i preporuke
    \item Status slučaja (otvoren, u toku, zatvoren)
    \item Ime i prezime veterinara / odgovorne osobe
\end{itemize}

\subsection*{Evidencija proizvodnje mlijeka po kravi}
Ovaj dokument sumira sve podatke o muži(izdajanju mlijeka) i količini proizvedenog mlijeka za svaku kravu. Većina polja se automatski popunjava preuzimanjem podataka sa sistema za mužu. \\

\noindent Elementi dokumenta su sljedeći:
\begin{itemize}
    \item Jedinstveni ID krave
    \item Datum muže
    \item Vrijeme početka i završetka muže
    \item Količina proizvedenog mlijeka (litri)
    \item Prosječan protok mlijeka
    \item Temperatura mlijeka na izlazu
    \item Broj muzne stanice / linije
    \item Način unosa podataka (automatski / ručni)
    \item Detektovana odstupanja (pad ili rast proizvodnje u odnosu na prosjek)
    \item Napomene (nepotpuna muža, tehnički kvar i sl.)
\end{itemize}

\subsection*{Evidencija uslova u štali (senzorska očitanja)}
Ovaj dokument predstavlja zapis očitanja IoT senzora koji prate okolišne uslove u štali i eventualna alarmantna stanja. \\

\noindent Elementi dokumenta su sljedeći:
\begin{itemize}
    \item Datum i vrijeme očitanja
    \item Oznaka štale / zone
    \item Temperatura zraka
    \item Relativna vlažnost
    \item Koncentracija amonijaka (NH$_3$) i CO$_2$
    \item Oznaka senzora ili uređaja
    \item Informacija o eventualnom prekoračenju pragova (normalno / upozorenje / kritično)
    \item Ime i prezime osobe koja je reagovala na alarm (ako postoji intervencija)
\end{itemize}

\newpage

\subsection*{Radni nalog / zadatak osoblju farme}
Na osnovu detekcija sistema ili plana rada generiše se radni nalog koji se dodjeljuje konkretnom članu osoblja. \\

\noindent Elementi dokumenta su sljedeći:
\begin{itemize}
    \item Jedinstveni ID zadatka
    \item Naziv zadatka (npr. „Veterinarski pregled krave 1234“)
    \item Detaljan opis posla
    \item Povezana krava / grupa krava ili zona štale
    \item Prioritet zadatka (normalan, visok, hitan)
    \item Osoba koja je kreirala nalog (menadžer, veterinar, sistem)
    \item Dodijeljeni izvršilac / tim
    \item Datum i vrijeme kreiranja zadatka
    \item Rok izvršenja
    \item Stvarno vrijeme početka i završetka
    \item Evidencija utrošenih resursa (lijekovi, hrana, materijal, radni sati)
    \item Status zadatka (planiran, u toku, završen, odgođen)
    \item Napomene i komentari nakon izvršenja
\end{itemize}

\subsection*{Hitno upozorenje o stanju krava ili uslova na farmi}
Kada sistem detektuje kritično odstupanje (bolest, nagli pad proizvodnje, loše uslove u štali), generiše se dokument – obavijest o hitnom upozorenju. \\

\noindent Elementi dokumenta su sljedeći:
\begin{itemize}
    \item Tip upozorenja (zdravstveni incident, pad proizvodnje, okolišni problem)
    \item Nivo kritičnosti (informativno, visoko, HITNO)
    \item ID krave ili zone na koju se upozorenje odnosi
    \item Kratki opis problema i razlog aktiviranja (AI nalaz, senzorsko očitanje)
    \item Datum i vrijeme detekcije
    \item Kanali preko kojih je upozorenje poslano (SMS, e-mail, mobilna aplikacija)
    \item Osoba koja je preuzela slučaj
    \item Vrijeme prve reakcije
    \item Preduzete mjere (npr. izolacija životinje)
    \item Trenutni status upozorenja (u obradi, riješeno, lažni alarm)
\end{itemize}

\subsection*{Evidencija zaposlenih i korisnika sistema}
Ovaj dokument sadrži osnovne podatke o zaposlenima na farmi i njihovim korisničkim nalozima u informacionom sistemu. \\

\noindent Elementi dokumenta su sljedeći:
\begin{itemize}
    \item Ime i prezime zaposlenog
    \item Radno mjesto / uloga (administrator, menadžer, veterinar, osoblje farme, finansijski analitičar)
    \item Korisničko ime u sistemu
    \item Kontakt podaci (telefon, e-mail)
    \item Datum zaposlenja
    \item Odgovorni odjel ili smjena
    \item Nivo ovlaštenja u sistemu
    \item Status naloga (aktivan, blokiran, arhiviran)
    \item Napomene (posebne kompetencije, ovlaštenja)
\end{itemize}

\subsection*{Godišnji izvještaj o poslovanju farme}
Godišnji izvještaj sumira ključne pokazatelje rada pametne farme i koristi se za strateško planiranje i izvještavanje prema vlasnicima i partnerima. \\

\noindent Elementi dokumenta su sljedeći:
\begin{itemize}
    \item Ukupan broj grla i broj muznih krava
    \item Ukupno proizvedena količina mlijeka u godini
    \item Prosječna proizvodnja mlijeka po kravi
    \item Broj telenja, slučajeva bolesti i smrti
    \item Ukupni troškovi liječenja i veterinarskih intervencija
    \item Ukupni troškovi hrane i ostalih potrošnih resursa
    \item Ostvareni prihodi od mlijeka i drugih proizvoda
    \item Ostvareni profit i poređenje sa prethodnom godinom
    \item Ključni uočeni trendovi (na osnovu AI analize) i preporuke za naredni period
    \item Vrijeme i mjesto podnošenja izvještaja
    \item Ime i prezime odgovorne osobe i potpis
\end{itemize}


\clearpage
\vspace*{\fill}

\begin{center}
    \includegraphics[width=\textwidth]{images/analizaDokumenata_Dijagram.png}

    \vspace{1em}
    {\large \textbf{Dijagram analize dokumenata}}
\end{center}

\vspace*{\fill}
\clearpage


\newpage
\section{DEFINICIJA ZAHTJEVA}
\label{sec:definicija-zahtjeva}


Ova sekcija definiše skup funkcionalnih i nefunkcionalnih zahtjeva za IS pametne farme krava. Svaki zahtjev ima jedinstveni identifikator, prioritet (MoSCoW: M/S/C) i mjerljive kriterije prihvatanja.


\newcolumntype{L}[1]{>{\raggedright\arraybackslash}p{#1}}
\newcolumntype{C}[1]{>{\centering\arraybackslash}p{#1}}

\subsection{Funkcionalni zahtjevi}

\renewcommand{\arraystretch}{1.25}
\setlength{\tabcolsep}{5pt}
\begin{longtable}{|L{1.6cm}|L{8.0cm}|C{0.9cm}|L{6.2cm}|}
\hline
\rowcolor[gray]{0.9}\textbf{ID} & \textbf{Opis} & \textbf{Prio} & \textbf{Kriteriji prihvatanja} \\
\hline
FR-01 & Registracija i evidencija krava (osnovni biom., porijeklo, datum rođenja/dolaska, status). & M &
\begin{minipage}[t]{6.2cm}\vspace{2pt}
\begin{itemize}\setlength\itemsep{2pt}\setlength\parskip{0pt}\setlength\parsep{0pt}
    \item Unos traje $\leq$ 3 min po kravi
    \item Polja validirana; duplikat se sprječava
    \item Generiše se jedinstveni ID i QR kod
\end{itemize}\vspace{2pt}
\end{minipage} \\
\hline
FR-02 & Neinvazivno \textbf{prepoznavanje jedinki} kamerom i AI (fleke). & M &
\begin{minipage}[t]{6.2cm}\vspace{2pt}
\begin{itemize}\setlength\itemsep{2pt}\setlength\parskip{0pt}\setlength\parsep{0pt}
    \item Tačnost identifikacije $\geq$ 95\% (test set)
    \item Vrijeme prepoznavanja $\leq$ 2 s po jedinki
    \item Ručna potvrda/override uz log
\end{itemize}\vspace{2pt}
\end{minipage} \\
\hline
FR-03 & Automatska \textbf{evidencija proizvodnje mlijeka} po muži (integracija sa aparatima/senzorima). & M &
\begin{minipage}[t]{6.2cm}\vspace{2pt}
\begin{itemize}\setlength\itemsep{2pt}
    \item Mjerenje greške $\leq$ 2\%
    \item Bilježi svaku mužu (vrijeme, trajanje, količina)
    \item Podržan manualni unos sa oznakom
\end{itemize}\vspace{2pt}
\end{minipage} \\
\hline
FR-04 & \textbf{Trendovi i prognoze proizvodnje} (dnevno/sedmično/mjesečno po kravi i stadu). & S &
\begin{minipage}[t]{6.2cm}\vspace{2pt}
\begin{itemize}\setlength\itemsep{2pt}
    \item Dashboard prikazuje trend i YoY/MoM poređenja
    \item Export kao PDF/CSV u jednom kliku
\end{itemize}\vspace{2pt}
\end{minipage} \\
\hline
FR-05 & \textbf{Monitoring zdravlja i ponašanja} (video + IoT + analitika). & M &
\begin{minipage}[t]{6.2cm}\vspace{2pt}
\begin{itemize}\setlength\itemsep{2pt}
    \item Detekcija anomalije u $<12$ h za $\geq$ 90\% slučajeva
    \item Stopa lažnih alarma $<5$\%
    \item Svaki slučaj vezan uz profil krave i historiju
\end{itemize}\vspace{2pt}
\end{minipage} \\
\hline
FR-06 & \textbf{Integracija IoT senzora} (temperatura, vlaga, NH$_3$, CO$_2$). & M &
\begin{minipage}[t]{6.2cm}\vspace{2pt}
\begin{itemize}\setlength\itemsep{2pt}
    \item Očitavanje najmanje svaka 5 min
    \item Gubitak paketa $<1$\% dnevno
    \item Pragovi i alarmne politike konfigurabilni
\end{itemize}\vspace{2pt}
\end{minipage} \\
\hline
FR-07 & \textbf{Upozorenja i notifikacije} (push/e-mail/SMS) po prioritetu. & M &
\begin{minipage}[t]{6.2cm}\vspace{2pt}
\begin{itemize}\setlength\itemsep{2pt}
    \item Slanje u $\leq$ 60 s od događaja
    \item Prioriteti: inform., visoko, HITNO
    \item Personalizacija kanala po korisniku
\end{itemize}\vspace{2pt}
\end{minipage} \\
\hline
FR-08 & \textbf{Upravljanje zadacima} (kreiranje, dodjela, praćenje, statusi). & S &
\begin{minipage}[t]{6.2cm}\vspace{2pt}
\begin{itemize}\setlength\itemsep{2pt}
    \item Dodjela zadatka u $\leq$ 3 min
    \item Evidencija vremena i resursa
    \item Veza sa kravom/zonama/senzorima
\end{itemize}\vspace{2pt}
\end{minipage} \\
\hline
FR-09 & \textbf{Generisanje izvještaja} (proizvodnja, zdravlje, reprodukcija, resursi). & M &
\begin{minipage}[t]{6.2cm}\vspace{2pt}
\begin{itemize}\setlength\itemsep{2pt}
    \item Standardni izvještaji $<30$ s
    \item Export: PDF, Excel/CSV
    \item Zakazivanje (dnevno/sedmično/mjesečno)
\end{itemize}\vspace{2pt}
\end{minipage} \\
\hline
FR-10 & \textbf{Upravljanje korisnicima i ulogama} (RBAC). & M &
\begin{minipage}[t]{6.2cm}\vspace{2pt}
\begin{itemize}\setlength\itemsep{2pt}
    \item Uloge: Admin, Menadžer, Veterinar, Osoblje, ReadOnly
    \item Audit svakog pristupa i promjene
    \item 2FA obavezno za administratore
\end{itemize}\vspace{2pt}
\end{minipage} \\
\hline
FR-11 & \textbf{Offline rad i sinhronizacija} za terenske uređaje. & S &
\begin{minipage}[t]{6.2cm}\vspace{2pt}
\begin{itemize}\setlength\itemsep{2pt}
    \item Unosi dostupni offline
    \item Auto-sinhronizacija po povratku mreže
    \item Rješavanje konflikata kroz pregled razlika
\end{itemize}\vspace{2pt}
\end{minipage} \\
\hline
FR-12 & \textbf{Konfiguracija sistema} (pragovi, zone, uređaji, jedinice mjere). & S &
\begin{minipage}[t]{6.2cm}\vspace{2pt}
\begin{itemize}\setlength\itemsep{2pt}
    \item Promjene bez restarta sistema
    \item Verziona kontrola konfiguracija
\end{itemize}\vspace{2pt}
\end{minipage} \\
\hline
FR-13 & \textbf{Integracije} sa vanjskim sistemima (mljekara ERP, veterinarski IS). & C &
\begin{minipage}[t]{6.2cm}\vspace{2pt}
\begin{itemize}\setlength\itemsep{2pt}
    \item REST API (JSON) sa dokumentacijom
    \item Webhook za događaje (alarm, završena muža)
\end{itemize}\vspace{2pt}
\end{minipage} \\
\hline
FR-14 & \textbf{Kvalitet mlijeka} (npr. temperatura na izlazu, provjera odstupanja). & S &
\begin{minipage}[t]{6.2cm}\vspace{2pt}
\begin{itemize}\setlength\itemsep{2pt}
    \item Pragovi po standardu korisnika
    \item Automatski alarm i blokada muže za kritične slučajeve
\end{itemize}\vspace{2pt}
\end{minipage} \\
\hline
FR-15 & \textbf{Audit i revizija} (nepromjenjivi logovi aktivnosti). & M &
\begin{minipage}[t]{6.2cm}\vspace{2pt}
\begin{itemize}\setlength\itemsep{2pt}
    \item Čuvanje logova min. 24 mjeseca
    \item Pretraga po korisniku, događaju i periodu
\end{itemize}\vspace{2pt}
\end{minipage} \\
\hline
\end{longtable}

\subsection{Nefunkcionalni zahtjevi}

\begin{longtable}{|L{1.8cm}|L{3.5cm}|L{8.8cm}|L{2.6cm}|}
\hline
\rowcolor[gray]{0.9}\textbf{ID} & \textbf{Kategorija} & \textbf{Zahtjev (mjera)} & \textbf{Provjera} \\
\hline
NFR-01 & Performanse (UI) & Latencija tipičnih operacija $<2$ s; kompleksne (analitika/izvještaj) $<5$ s. & Test mjerenja \\
\hline
NFR-02 & Dostupnost & Uptime sistema $\geq$ 99\% mjesečno; planirani prekidi najavljeni 48h ranije. & Monitoring \\
\hline
NFR-03 & Skalabilnost & Horizontalno skalirati na $\geq$ 50 simultanih korisnika i 1{,}500+ krava bez degradacije. & Load test \\
\hline
NFR-04 & Pouzdanost senzora & Pouzdanost isporuke IoT podataka $\geq$ 99\%; re-try i buffer na rubu (edge). & Test/Sim. \\
\hline
NFR-05 & Sigurnost – prijenos & Šifrovanje u prijenosu TLS~1.2+; HSTS uključen. & Security scan \\
\hline
NFR-06 & Sigurnost – pohrana & Podaci u mirovanju: AES-256; lozinke hashirane (Argon2 ili bcrypt $>$ 10 rounds). & Code review \\
\hline
NFR-07 & Pristup/Autentikacija & 2FA za administratore; password politika $\geq$ 12 znakova; session timeout 30 min. & Pen-test \\
\hline
NFR-08 & Privatnost i usklađenost & Usklađeno sa GDPR gdje primjenjivo; \emph{privacy by design}. Video snimci čuvaju se 30 dana. & Audit \\
\hline
NFR-09 & Backup i oporavak & Dnevni full + inkrementalni backup; RPO $\leq$ 15 min, RTO $\leq$ 4 h. & DR drill \\
\hline
NFR-10 & Upotrebljivost & 90\% tipičnih zadataka izvedivo u $\leq$ 3 klika; onboarding tutorijali. & UX test \\
\hline
NFR-11 & Pristupačnost & Web klijent ispunjava WCAG~2.1~AA (kontrast, tastatura, alt opisi). & A11y review \\
\hline
NFR-12 & Kompatibilnost & Podržani preglednici: zadnje 2 verzije Chrome/Edge/Firefox i Safari; responzivno 360–1920 px. & Cross-browser \\
\hline
NFR-13 & Lokalizacija & Interfejs najmanje na \emph{bs/hr/sr} i \emph{en}; lokalizovani datumi, jedinice i valute. & Pregled \\
\hline
NFR-14 & Observability & Centralizovani logovi/metrici/alerti (CPU, latencija, greške); korelacija trace-ID. & Ops test \\
\hline
NFR-15 & Kvalitet AI modela & ID tačnost $\geq$ 95\%; zdravstvena detekcija $\geq$ 90\%; lažni alarmi $<5$\%; periodična re-obuka. & Eval set \\
\hline
NFR-16 & Integrabilnost & Javne REST API rute sa OpenAPI specifikacijom; rate-limit i API ključ / OAuth2. & API test \\
\hline
NFR-17 & Održavanje & Pokrivenost testovima $\geq$ 70\% kritičnog koda; CI/CD sa statičkom analizom. & CI izvještaj \\
\hline
NFR-18 & Upravljanje podacima & Retencija poslovnih podataka konfigurabilna; audit trail $\geq$ 24 mjeseca. & Audit \\
\hline
NFR-19 & Offline režim & Klijent kešira zapise i sinhronizuje u $\leq$ 15 min nakon povratka mreže. & Field test \\
\hline
NFR-20 & Etika i dobrobit & AI odluke objašnjive (razlog alarma); sistem minimizira stres životinja (bez invazivnih tagova). & Review \\
\hline
\end{longtable}

\subsection{Ograničenja i pretpostavke}
\begin{itemize}
    \item Kamere pokrivaju $\geq$ 95\% prostora štale; osvjetljenje zadovoljava preporučene nivoe.
    \item Senzori posjeduju certifikate i redovnu kalibraciju prema planu održavanja.
    \item Mrežna infrastruktura omogućava RTSP/MQTT sa stabilnim protokom i QoS politikama.
\end{itemize}

\newpage
\section{SLUČAJEVI UPOTREBE}

% US-1: Evidencija krava
\begin{table}[H]
\centering
\small
\begin{tabular}{|p{3.5cm}|p{11cm}|}
\hline
\rowcolor[HTML]{B4C7E7}
\multicolumn{2}{|c|}{\textbf{US-1: Evidencija krava}} \\
\hline
\textbf{Proces:} & Registracija i evidencija novih krava. \\
\hline
\textbf{ID:} & US-1 \\
\hline
\textbf{Prioritet:} & Visoki \\
\hline
\textbf{Učestalost:} & 5–10 puta mjesečno \\
\hline
\textbf{Učesnik:} & Osoblje farme, Administrator \\
\hline
\textbf{Opis:} & Unos osnovnih, biometrijskih i proizvodnih podataka o novoj kravi. Sistem automatski generiše jedinstveni ID na osnovu fleka. \\
\hline
\textbf{Trigger:} & Dolazak nove krave (kupovina, rođenje), potreba za ažuriranjem, migracija iz starog sistema. \\
\hline
\textbf{Preduvjeti:} & Autentifikovan korisnik sa dozvolom; sistem i baza dostupni. \\
\hline
\textbf{Normalan tok:} &
\begin{enumerate}[leftmargin=*, nosep]
    \item Otvaranje modula za evidenciju
    \item Prikaz formulara
    \item Unos osnovnih podataka (rođenje, rasa, porijeklo, matični broj, težina, reproduktivni status)
    \item Snimanje fotografija iz više uglova
    \item AI analizira raspored fleka i kreira biometrijski profil
    \item Validacija podataka
    \item Generisanje jedinstvenog ID-a i spremanje u bazu
    \item Prikaz potvrde registracije
    \item Automatsko kreiranje profila proizvodnje i zdravlja
    \item Generisanje QR koda
\end{enumerate} \\
\hline
\textbf{Alternativni tok:} &
\textit{A1 – Nevalidni podaci:} Sistem označava greške, korisnik ispravlja. \newline
\textit{A2 – Duplikat:} Detekcija duplikata; korisnik može ažurirati, nastaviti ili otkazati. \newline
\textit{A3 – AI ne može prepoznati:} Traži dodatne fotografije; ako ne uspije, koristi se manuelna identifikacija. \\
\hline
\textbf{Tok izuzetaka:} &
\textit{E1 – Greška baze:} Poništavanje transakcije, prikaz poruke. \newline
\textit{E2 – Gubitak konekcije:} Lokalno privremeno čuvanje. \newline
\textit{E3 – Neispravna kamera:} Korištenje mobilnog uređaja. \newline
\textit{E4 – Sistemski pad:} Automatski recovery. \\
\hline
\textbf{Rezultat:} & Krava registrovana, generisan ID, kreiran profil, AI model ažuriran, poslata notifikacija administratoru. \\
\hline
\textbf{Ulazi/Izlazi:} &
\textit{Ulazi:} Osnovni podaci, fotografije krave. \newline
\textit{Izlazi:} ID, biometrijski profil, QR kod, potvrda, notifikacija. \\
\hline
\end{tabular}
\caption{Slučaj upotrebe US-1: Evidencija krava}
\end{table}


\pagebreak

% US-2: Evidencija proizvodnje mlijeka
\begin{table}[H]
\centering
\small
\begin{tabular}{|p{3.5cm}|p{11cm}|}
\hline
\rowcolor[HTML]{B4C7E7}
\multicolumn{2}{|c|}{\textbf{US-2: Evidencija proizvodnje mlijeka}} \\
\hline
\textbf{Proces:} & Automatsko i manuelno praćenje proizvodnje mlijeka po kravi. \\
\hline
\textbf{ID:} & US-2 \\
\hline
\textbf{Prioritet:} & Kritični \\
\hline
\textbf{Učestalost:} & Vrlo često (1600–2400 muži dnevno). \\
\hline
\textbf{Učesnik:} & Osoblje farme, Administrator, Sistem za mužu \\
\hline
\textbf{Opis:} & Sistem automatski bilježi količinu mlijeka po muži, analizira trendove, detektuje odstupanja i generiše upozorenja. Podaci služe za optimizaciju ishrane i rano otkrivanje zdravstvenih problema. \\
\hline
\textbf{Trigger:} & Završena muža, signal iz sistema za mužu, manuelni unos, zahtjev za izvještaj. \\
\hline
\textbf{Preduvjeti:} & Krava je registrovana; sistem za mužu kalibrisan; senzori aktivni; AI modul za prepoznavanje krave funkcionalan. \\
\hline
\textbf{Normalan tok:} & 
\begin{enumerate}[leftmargin=*, nosep]
    \item Krava ulazi u boks
    \item AI prepoznaje kravu
    \item Učitava se profil i historija
    \item Započinje muža
    \item Senzori mjere protok i količinu
    \item Sistem bilježi količinu, trajanje, protok, kvalitet, vrijeme
    \item AI analizira podatke u realnom vremenu
    \item Poređenje sa historijom
    \item Ako je normalno → automatsko spremanje
    \item Ažuriraju se preporuke za ishranu
    \item Prikazuje se potvrda muže
    \item Ažuriranje dnevnih/sedmičnih/mjesečnih izvještaja
\end{enumerate} \\
\hline
\textbf{Alternativni tok:} & 
\textit{A1 – Pad proizvodnje:} Pad >15\%, alarm, notifikacija veterinaru/osoblju, označavanje krave, prijedlog uzroka. \newline
\textit{A2 – Povećana proizvodnja:} Porast >20\%, pozitivna notifikacija, prilagodba ishrane. \newline
\textit{A3 – Manuelni unos:} Osoblje bira kravu i unosi količinu; unos se označava kao manuelan. \newline
\textit{A4 – Neprepoznata krava:} AI nesiguran → prikazuje top 3; osoblje bira; evidentira se manual override. \\
\hline
\textbf{Izuzeci:} & 
\textit{E1 – Senzor ne radi:} Manuelni unos, oznaka za servis. \newline
\textit{E2 – Kamera ne radi:} RFID ili manuelni unos. \newline
\textit{E3 – Prekid napajanja:} Automatski recovery, procjena količine. \newline
\textit{E4 – Greška baze:} Lokalno čuvanje, kasnija sinkronizacija. \newline
\textit{E5 – Loš kvalitet mlijeka:} Zaustavljanje muže, obavještavanje veterinara. \\
\hline
\textbf{Rezultat:} & Podaci spremljeni, statistika ažurirana, generisani trendovi i prognoze, kreirani alarmi, podaci integrisani s modulom za ishranu. \\
\hline
\textbf{Ulazi/Izlazi:} & 
\textit{Ulazi:} Identifikacija krave; senzori (količina, protok, temperatura). \newline
\textit{Izlazi:} Sačuvani podaci, trendovi, statistika, upozorenja, preporuke ishrane, izvještaji. \\
\hline
\end{tabular}
\caption{Slučaj upotrebe US-2: Evidencija proizvodnje mlijeka}
\end{table}


\pagebreak

% US-3: Praćenje zdravstvenog stanja
\begin{table}[H]
\centering
\small
\begin{tabular}{|p{3.5cm}|p{11cm}|}
\hline
\rowcolor[HTML]{B4C7E7}
\multicolumn{2}{|c|}{\textbf{US-3: Praćenje zdravstvenog stanja}} \\
\hline
\textbf{Proces:} & Kontinuirani monitoring zdravstvenog stanja i ponašanja krava. \\
\hline
\textbf{ID:} & US-3 \\
\hline
\textbf{Prioritet:} & Kritični \\
\hline
\textbf{Učestalost:} & 24/7 (10–15 upozorenja dnevno) \\
\hline
\textbf{Učesnik:} & Veterinar, Administrator \\
\hline
\textbf{Opis:} & AI, video-analiza i IoT senzori kontinuirano prate vitale i ponašanje. Sistem automatski detektuje odstupanja i omogućava ranu intervenciju. \\
\hline
\textbf{Trigger:} & Kontinuirani monitoring; anomalije; pad mlijeka; promjena temperature; ručna provjera; zakazani pregled. \\
\hline
\textbf{Preduvjeti:} & Aktivne kamere i senzori; obučen AI model; krava ima profil; uspostavljena bazna linija; veterinar ima pristup. \\
\hline
\textbf{Normalan tok:} & 
\begin{enumerate}[leftmargin=*, nosep]
    \item Prikupljanje podataka (video, temperatura, aktivnost, proizvodnja, interakcije)
    \item AI analiza svakih 30 min
    \item Poređenje s baznom linijom, stadom i historijom
    \item Ako je sve normalno → status „zdrava“, arhiviranje
    \item Ako postoji anomalija → povećana analiza, dodatni podaci, klasifikacija, alarm
    \item Slanje notifikacija veterinaru/osoblju (ID krave, simptomi, video, dijagnoza, preporuke)
    \item Veterinar verifikuje dijagnozu i propisuje tretman
    \item Osoblje izvršava tretman
    \item Sistem prati efikasnost terapije
    \item Povratak u normalan status
    \item Dokumentovanje slučaja za AI obuku
\end{enumerate} \\
\hline
\textbf{Alternativni tok:} & 
\textit{A1 – Hitno:} Kritični simptomi, HITNO upozorenje, automatski poziv veterinara, SMS, fokusirani video nadzor. \newline
\textit{A2 – Lažni alarm:} Veterinar potvrđuje lažni alarm; AI dopunska obuka. \newline
\textit{A3 – Estrus:} Detekcija estrusa, označavanje krave, notifikacija inseminatoru, optimalno vrijeme. \newline
\textit{A4 – Grupna anomalija:} Više krava pokazuje simptome → alarm za epidemiju, izolacija, obavještenje institucijama. \\
\hline
\textbf{Izuzeci:} & 
\textit{E1 – Kvar senzora/kamere:} Redundantni uređaji, servis. \newline
\textit{E2 – Nejasni AI rezultati:} Ručna provjera. \newline
\textit{E3 – Nedostupan veterinar:} Backup veterinar, eskalacija. \newline
\textit{E4 – Prekid napajanja:} UPS, mobilni pristup. \newline
\textit{E5 – Greška baze:} Lokalni log, hitna sinkronizacija. \\
\hline
\textbf{Rezultat:} & Ažurirano zdravstveno stanje, kreirane notifikacije, dokumentovana dijagnoza i tretman, ažuriran karton, generisani izvještaji. \\
\hline
\textbf{Ulazi/Izlazi:} & 
\textit{Ulazi:} Video, senzori (temperatura, aktivnost), proizvodnja, historija. \newline
\textit{Izlazi:} Status, upozorenja, dijagnoze, preporuke, zdravstveni zapisi, izvještaji. \\
\hline
\end{tabular}
\caption{Slučaj upotrebe US-3: Praćenje zdravstvenog stanja}
\end{table}


\pagebreak

% US-4: Upravljanje zadacima
\begin{table}[H]
\centering
\small
\begin{tabular}{|>{\raggedright}p{3.5cm}|>{\raggedright\arraybackslash}p{11cm}|}
\hline
\rowcolor[HTML]{B4C7E7}
\multicolumn{2}{|c|}{\textbf{US-4: Upravljanje zadacima}} \\
\hline
\textbf{Proces:} & Kreiranje, dodjeljivanje i praćenje svakodnevnih i periodičnih zadataka na farmi. \\
\hline
\textbf{ID:} & US-4 \\
\hline
\textbf{Prioritet:} & Srednji \\
\hline
\textbf{Učestalost:} & Često (20–50 novih dnevno, 100+ rutinskih mjesečno) \\
\hline
\textbf{Učesnik:} & Menadžer, Administrator, Osoblje farme \\
\hline
\textbf{Opis:} & Sistem centralizuje upravljanje zadacima, automatski generiše rutinske aktivnosti, prati izvršenje, bilježi resurse i integriše se s modulima zdravlja i proizvodnje. \\
\hline
\textbf{Trigger:} & Početak dana, kreiranje zadatka, sistemska upozorenja (vakcinacija, pregled), izmjena rasporeda, završetak zadatka. \\
\hline
\textbf{Preduvjeti:} & Aktivni korisnički pristup, registrovano osoblje, definisani tipovi/prioriteti zadataka, aktivan kalendar farme. \\
\hline
\textbf{Normalan tok:} & 
\begin{enumerate}[leftmargin=*, nosep]
    \item Menadžer otvara modul zadataka.
    \item Sistem prikazuje kalendar, aktivne zadatke, status i dostupnost osoblja.
    \item Menadžer kreira zadatak: tip, prioritet, opis, trajanje, resursi.
    \item Sistem predlaže dostupno osoblje prema kvalifikacijama, opterećenju, lokaciji i historiji.
    \item Menadžer dodjeljuje zadatak.
    \item Sistem validira dostupnost i konflikte.
    \item Sistem šalje notifikacije (push, email, dashboard).
    \item Osoblje prima i potvrđuje zadatak.
    \item Tokom izvršenja ažurira status, dodaje bilješke/fotografije, vrijeme i resurse.
    \item Po završetku označava zadatak kao završen.
    \item Menadžer odobrava završetak.
    \item Sistem arhivira zadatak i ažurira statistiku.
\end{enumerate} \\
\hline
\textbf{Alternativni tokovi:} &
A1 – Automatski rutinski zadaci (hranjenje, muža, čišćenje).  
A2 – Zdravstveni zadatak: AI kreira hitan zadatak.  
A3 – Nedostupno osoblje: sistem predlaže zamjenu ili odgodu.  
A4 – Problem u izvršenju: osoblje prijavljuje, menadžer dobija obavijest.  
A5 – Odbijen završetak: menadžer vraća zadatak na doradu. \\
\hline
\textbf{Izuzeci:} & 
E1 – Sistemski kvar: ručno kreiranje i obavještavanje.  
E2 – Bez konekcije: lokalno čuvanje i kasnija sinhronizacija.  
E3 – Nema dostupnog osoblja: alarm i eskalacija.  
E4 – Konflikt prioriteta: automatsko rangiranje, menadžer odlučuje. \\
\hline
\textbf{Rezultat:} & Zadatak kreiran/dodijeljen, poslane notifikacije, ažuriran kalendar i dashboard, resursi rezervisani, evidencija izvršenja sačuvana. \\
\hline
\textbf{Ulazi / Izlazi:} & 
Ulazi: tip, prioritet, opis, trajanje, resursi, dodijeljeno osoblje.  
Izlazi: kreiran zadatak, notifikacije, ažuriran kalendar, evidencija izvršenja, statistika produktivnosti. \\
\hline
\end{tabular}
\caption{Slučaj upotrebe US-4: Upravljanje zadacima}
\end{table}


\pagebreak

% US-5: Generisanje izvještaja
\begin{table}[H]
\centering
\small
\begin{tabular}{|>{\raggedright}p{3.5cm}|>{\raggedright\arraybackslash}p{11cm}|}
\hline
\rowcolor[HTML]{B4C7E7}
\multicolumn{2}{|c|}{\textbf{US-5: Generisanje izvještaja}} \\
\hline
\textbf{Proces:} & Kreiranje analitičkih izvještaja o proizvodnji, zdravlju i performansama farme. \\
\hline
\textbf{ID:} & US-5 \\
\hline
\textbf{Prioritet:} & Visoki \\
\hline
\textbf{Učestalost:} & Često (30–60 izvještaja mjesečno) \\
\hline
\textbf{Učesnik:} & Menadžer, Veterinar, Finansijski analitičar, Administrator \\
\hline
\textbf{Opis:} & Sistem generiše sveobuhvatne izvještaje uz automatsku agregaciju podataka, vizualizaciju trendova, prediktivnu analitiku i izvoz u više formata. \\
\hline
\textbf{Trigger:} & Manuelni zahtjev, automatski periodični izvještaji (dnevni/sedmični/mjesečni), kraj perioda, audit zahtjev, integracije. \\
\hline
\textbf{Preduvjeti:} & Korisnik ima prava pristupa; podaci dostupni; definisan period; analitički modul operativan. \\
\hline
\textbf{Normalan tok:} & 
\begin{enumerate}[leftmargin=*, nosep]
    \item Korisnik otvara modul izvještaja.
    \item Sistem prikazuje tipove (proizvodnja, zdravlje, finansije, reprodukcija, resursi, efikasnost, kvalitet, komparacije).
    \item Korisnik bira tip izvještaja.
    \item Sistem nudi parametre: period, detaljnost, filtri, format, vizualizacije.
    \item Korisnik postavlja parametre i pokreće generisanje.
    \item Sistem provjerava dostupne podatke.
    \item Sistem vrši obradu: agregacija, statistike, vizualizacije, AI analiza trendova i predikcija.
    \item Generiše se izvještaj (summary, KPI-ovi, analize, vizualizacije, komparacije, prognoze, problemi, preporuke).
    \item Izvještaj se prikazuje u web interfejsu.
    \item Korisnik vrši interaktivnu analizu (drill-down, prilagođavanje grafova, komentari).
    \item Izvještaj se eksportuje.
    \item Sistem arhivira izvještaj.
    \item Kod periodičnih izvještaja zakazuje se sljedeći termin. 
\end{enumerate} \\
\hline
\textbf{Alternativni tokovi:} & 
A1 – Nedostajući podaci (upozorenje, parcijalni izvještaj, odgoda, procjene).  
A2 – Custom izvještaj (builder: izvori, metrike, grafovi, layout, template).  
A3 – Automatsko zakazivanje (frekvencija, vrijeme, primaoci, format).  
A4 – Komparativna analiza (raniji period, prošla godina, industrijski prosjek).  
A5 – Sporo generisanje ( $ >30s =>$ progress bar, notifikacija). \\
\hline
\textbf{Izuzeci:} & 
E1 – Greška u bazi (recovery, log, notifikacija).  
E2 – Timeout (predlaže smanjenje perioda/detalja).  
E3 – Neispravni parametri (validacija, sugestije).  
E4 – Nedovoljno prostora (čišćenje izvještaja $>6$ mjeseci).  
E5 – Greška pri eksportu (alternativni format ili prikaz). \\
\hline
\textbf{Rezultat:} & Izvještaj generisan, vizualizovan, eksportovan, arhiviran; kreirane preporuke. \\
\hline
\textbf{Ulazi / Izlazi:} & 
Ulazi: tip izvještaja, period, filtri, detaljnost, format.  
Izlazi: izvještaj, vizualizacije, KPI-ovi, trendovi, prognoze, preporuke, eksportovan fajl. \\
\hline
\end{tabular}
\caption{Slučaj upotrebe US-5: Generisanje izvještaja}
\end{table}


\pagebreak

% US-6: Upravljanje korisnicima
\begin{table}[H]
\centering
\small
\begin{tabular}{|>{\raggedright}p{3.5cm}|>{\raggedright\arraybackslash}p{11cm}|}
\hline
\rowcolor[HTML]{B4C7E7}
\multicolumn{2}{|c|}{\textbf{US-6: Upravljanje korisnicima}} \\
\hline
\textbf{Proces:} & Administracija korisničkih naloga, uloga i pristupnih prava. \\
\hline
\textbf{ID:} & US-6 \\
\hline
\textbf{Prioritet:} & Visok (sigurnosno kritičan) \\
\hline
\textbf{Učestalost:} & Povremeno (2–5 novih mjesečno, 10–20 izmjena, 1–3 deaktivacije, kvartalne revizije). \\
\hline
\textbf{Učesnik:} & Administrator \\
\hline
\textbf{Opis:} & Sistem omogućava kreiranje, uređivanje, deaktivaciju i brisanje naloga uz RBAC model dozvola i detaljan audit log. \\
\hline
\textbf{Trigger:} & Novi zaposlenik, promjena uloge, odlazak zaposlenika, zahtjev za reset lozinke, sigurnosni incident, periodična revizija. \\
\hline
\textbf{Preduvjeti:} & Administrator je prijavljen, ima odgovarajuća prava; definisane uloge i politike pristupa; autentifikacijski sistem funkcionalan. \\
\hline
\textbf{Normalan tok:} & 
\begin{enumerate}[leftmargin=*, nosep]
    \item Administrator otvara modul za upravljanje korisnicima.
    \item Sistem prikazuje listu korisnika i njihov status.
    \item Administrator bira kreiranje korisnika.
    \item Sistem prikazuje formular (osnovni podaci, uloga, odjel).
    \item Administrator unosi podatke i odabire ulogu.
    \item Sistem validira email i lozinku.
    \item Administrator podešava specifična prava.
    \item Sistem kreira nalog i generiše privremenu lozinku.
    \item Sistem šalje email za aktivaciju.
    \item Korisnik aktivira nalog i postavlja lozinku.
    \item Sistem bilježi događaj u audit log.
\end{enumerate} \\
\hline
\textbf{Alternativni tokovi:} & 
A1 – Ažuriranje korisnika.  
A2 – Deaktivacija naloga.  
A3 – Resetovanje lozinke.  
A4 – Masovni import (CSV).  
A5 – Revizija pristupa (neaktivni $>90$ dana). \\
\hline
\textbf{Izuzeci:} & 
E1 – Email već postoji.  
E2 – Slaba lozinka.  
E3 – Neuspješno slanje emaila.  
E4 – Nedozvoljeno brisanje (aktivne obaveze).  
E5 – Sumnja na kompromitaciju → automatska blokada. \\
\hline
\textbf{Rezultat:} & Kreiran/ažuriran/deaktiviran nalog, ažurirane dozvole, poslane notifikacije i evidentiran audit log. \\
\hline
\textbf{Ulazi / Izlazi:} & 
Ulazi: korisnički podaci, uloga, dozvole.  
Izlazi: nalog, email obavijest, audit zapis, izvještaji. \\
\hline
\end{tabular}
\caption{Slučaj upotrebe US-6: Upravljanje korisnicima}
\end{table}

\newpage
\section{DIJAGRAM SLUČAJEVA UPOTREBE}

\begin{figure}[h]
    \centering
    \includegraphics[width=1\textwidth]{images/UseCase.jpg}
    \caption{Dijagram slučajeva upotrebe}
    \label{fig:my_image}
\end{figure}

\newpage
\section{ACTIVITY DIJAGRAMI}

\begin{figure}[H]
    \centering
    \includegraphics[width=1\textwidth]{images/Aktivnost1.drawio.png}
    \label{fig:my_image}
\end{figure}

\begin{figure}[H]
    \centering
    \includegraphics[width=1\textwidth]{images/AktivnostD2.drawio.png}
    \label{fig:my_image}
\end{figure}

\begin{figure}[H]
    \centering
    \includegraphics[width=0.9\textwidth]{images/ad-3.png}
    \label{fig:my_image}
\end{figure}

\begin{figure}[H]
    \centering
    \includegraphics[width=1\textwidth]{images/ad-4.png}
    \label{fig:my_image}
\end{figure}

\begin{figure}[H]
    \centering
    \includegraphics[width=0.8\textwidth]{images/ad_zana.png}
    \label{fig:my_image}
\end{figure}

\begin{figure}[H]
    \centering
    \includegraphics{images/AD6.png}
    \label{fig:my_image}
\end{figure}
\newpage
\section{ERD DIJAGRAM BAZE PODATAKA}

ERD (Entity--Relationship Diagram) prikazuje logički model baze podataka informacionog sistema
\textbf{„OIS -- Smart Cow Farm“}. Dijagram ilustruje glavne entitete, njihove atribute i veze koje
među njima postoje, a nastao je na osnovu prethodno definisanih funkcionalnih i nefunkcionalnih
zahtjeva sistema.

Na ERD dijagramu su posebno istaknuti sljedeći ključni entiteti:
\begin{itemize}
    \item \textbf{KRAVA} -- centralni entitet sistema, sadrži osnovne podatke o svakoj kravi
    (identifikacija, rasa, status, proizvodne karakteristike).
    \item \textbf{MUZA} -- bilježi pojedinačne muže i količinu proizvedenog mlijeka po kravi.
    \item \textbf{ZDRAVSTVENI\_SLUCAJ} -- evidencija zdravstvenih stanja i intervencija za svaku kravu.
    \item \textbf{TERAPIJA} i \textbf{TERAPIJA\_APLIKACIJA} -- definiraju propisane terapije i svaku
    pojedinačnu primjenu lijeka.
    \item \textbf{SENZOR} i \textbf{OCITANJA\_SENZORA} -- opisuju IoT senzore i historiju očitanja
    uslova okoline u štali.
    \item \textbf{ZADATAK} -- radni nalozi za osoblje farme (veterinarski pregledi, hranjenje, čišćenje, održavanje).
    \item \textbf{UPOZORENJE} -- sistemska obavještenja koja nastaju na osnovu AI analize, senzorskih očitanja
    ili pada proizvodnje.
    \item \textbf{KORISNIK} -- svi korisnici sistema (administrator, menadžer, veterinar, osoblje farme) sa svojim
    nalozima i ovlaštenjima.
\end{itemize}

ERD dijagram osigurava da je struktura baze podataka usklađena sa procesima na farmi, omogućava
normalizaciju podataka i olakšava kasniju implementaciju relacione baze podataka (npr. PostgreSQL ili MySQL).

\begin{figure}[H]
    \centering
    \includegraphics[width=\textwidth]{images/erd.png}
    \caption{ERD dijagram baze podataka sistema „OIS -- Smart Cow Farm"}
    \label{fig:erd_smart_cow_farm}
\end{figure}

\newpage
\section{DIZAJN ARHITEKTURE}

\subsection{Odabir arhitekture}

Za informacioni sistem „OIS -- Smart Cow Farm" odabrana je \textbf{mikroservisna arhitektura} deployirana na \textbf{K3s Kubernetes klasteru}. Ovaj izbor omogućava visoku dostupnost (High Availability), automatsko skaliranje servisa prema opterećenju, te jednostavno održavanje i nadogradnju pojedinačnih komponenti sistema bez prekida rada cjelokupne platforme.

\subsubsection{Obrazloženje izbora}

Mikroservisna arhitektura je odabrana iz sljedećih razloga:

\begin{itemize}
    \item \textbf{Nezavisno skaliranje} -- AI servisi za prepoznavanje krava i analizu zdravlja zahtijevaju GPU resurse i mogu se skalirati nezavisno od web aplikacije i API servisa
    \item \textbf{Izolacija grešaka} -- kvar jednog servisa (npr. video procesiranja) ne utječe na rad ostalih komponenti sistema
    \item \textbf{Tehnološka fleksibilnost} -- različiti servisi mogu koristiti najprikladnije tehnologije (Python za AI, Go/Rust za performanse, Node.js za real-time)
    \item \textbf{Kontinuirana isporuka} -- omogućava nezavisno deployiranje i ažuriranje servisa bez downtime-a
\end{itemize}

\subsubsection{K3s Kubernetes klaster}

K3s je odabran kao lightweight Kubernetes distribucija optimizirana za edge computing i IoT okruženja, što je idealno za farmu koja kombinira cloud i lokalne resurse. Klaster se sastoji od:

\begin{itemize}
    \item \textbf{3 control-plane čvora} -- osiguravaju visoku dostupnost upravljačke ravni
    \item \textbf{N worker čvorova} -- skaliraju se prema potrebama, uključujući GPU čvorove za AI workloade
    \item \textbf{Edge čvorovi na farmi} -- za nisku latenciju kod obrade video streama i senzorskih podataka
\end{itemize}

\subsection{Opis arhitekture sistema}

Sistem je organiziran u sljedeće logičke slojeve i servise:

\subsubsection{Prezentacijski sloj}

\begin{itemize}
    \item \textbf{Web aplikacija (SPA)} -- React/Vue.js aplikacija za menadžment i osoblje farme
    \item \textbf{Mobile-responsive dizajn} -- pristup putem mobilnih uređaja u štali
    \item \textbf{Real-time dashboard} -- prikaz živih podataka sa senzora i upozorenja
\end{itemize}

\subsubsection{API Gateway sloj}

\begin{itemize}
    \item \textbf{Traefik Ingress Controller} -- rutiranje zahtjeva, SSL terminacija, rate limiting
    \item \textbf{REST API} -- za CRUD operacije nad podacima, dokumentovano sa OpenAPI specifikacijom
    \item \textbf{WebSocket API} -- za real-time notifikacije i live streaming podataka
    \item \textbf{Rate limiting} -- zaštita od preopterećenja API-ja
    \item \textbf{Autentifikacija API-ja:} API ključevi za IoT uređaje, OAuth2 za eksterne integracije
\end{itemize}

\subsubsection{Aplikacijski sloj (Mikroservisi)}

\begin{enumerate}
    \item \textbf{Auth Service} -- autentifikacija (JWT tokeni), autorizacija, upravljanje korisnicima
    \item \textbf{Cow Management Service} -- CRUD operacije za krave, praćenje statusa
    \item \textbf{Milk Production Service} -- evidencija muža, statistika proizvodnje
    \item \textbf{Health Service} -- zdravstveni slučajevi, terapije, veterinarski pregledi
    \item \textbf{Task Service} -- upravljanje zadacima za osoblje
    \item \textbf{Alert Service} -- generisanje i distribucija upozorenja
    \item \textbf{Sensor Service} -- prijem i obrada IoT senzorskih podataka
    \item \textbf{AI Inference Service} -- prepoznavanje krava, analiza ponašanja, detekcija anomalija
    \item \textbf{Video Processing Service} -- obrada video streama sa kamera
    \item \textbf{Notification Service} -- email, SMS, push notifikacije
    \item \textbf{Report Service} -- generisanje izvještaja i analitika
\end{enumerate}

\subsubsection{Sloj podataka}

\begin{itemize}
    \item \textbf{PostgreSQL} -- primarna relaciona baza podataka (prema ERD dijagramu)
    \item \textbf{TimescaleDB} -- ekstenzija PostgreSQL-a za vremenske serije (senzorski podaci)
    \item \textbf{Redis} -- caching, session storage, message broker za real-time
    \item \textbf{MinIO} -- S3-kompatibilno objektno skladište za slike i video snimke
\end{itemize}

\subsubsection{AI/ML infrastruktura}

\begin{itemize}
    \item \textbf{NVIDIA Triton Inference Server} -- serving AI modela za prepoznavanje krava
    \item \textbf{Model Registry} -- verzioniranje i upravljanje ML modelima
    \item \textbf{Feature Store} -- skladištenje izračunatih feature-a za AI modele
\end{itemize}

\textbf{Ciljane metrike AI modela:}
\begin{itemize}
    \item Tačnost identifikacije krava: $\geq$95\%
    \item Tačnost detekcije zdravstvenih problema: $\geq$90\%
    \item Lažni alarmi: $<$5\%
    \item Vrijeme inference-a: $\leq$2 sekunde
    \item Periodična re-obuka modela (kvartalno ili pri padu performansi)
    \item Objašnjivost AI odluka (razlog alarma vidljiv korisniku)
    \item Neinvazivni pristup -- bez fizičkih tagova na životinjama
\end{itemize}

\subsubsection{IoT infrastruktura}

\begin{itemize}
    \item \textbf{MQTT Broker (EMQX)} -- prijem podataka sa IoT senzora, QoS politike za garantovanu isporuku
    \item \textbf{Edge Gateway} -- lokalna obrada na farmi, buffer pri gubitku konekcije, retry mehanizam
    \item \textbf{Pouzdanost isporuke:} $\geq$99\% senzorskih podataka uspješno isporučeno
    \item \textbf{Frekvencija očitanja:} Svakih 5 minuta, dnevni gubitak paketa $<$1\%
    \item \textbf{Automatska rekalibracija:} Senzori sa definisanim pragovima i politikom alarma
\end{itemize}

\subsection{Nefunkcionalni zahtjevi -- Arhitekturalna rješenja}

\subsubsection{Prenosivost i kompatibilnost}

\begin{itemize}
    \item \textbf{Docker kontejneri} -- svi servisi su pakovani u Docker kontejnere, osiguravajući konzistentno ponašanje na svim okruženjima (development, staging, production)
    \item \textbf{Kubernetes} -- standardizirana orkestracija omogućava migraciju između cloud providera
    \item \textbf{Infrastructure as Code (Terraform)} -- reproduktivna infrastruktura
    \item \textbf{Helm Charts} -- standardizirano deployiranje aplikacija
\end{itemize}

\subsubsection{Performanse i skalabilnost}

\begin{itemize}
    \item \textbf{Horizontal Pod Autoscaler (HPA)} -- automatsko skaliranje servisa prema CPU/memoriji
    \item \textbf{KEDA (Kubernetes Event-driven Autoscaling)} -- skaliranje prema custom metrikama (npr. broj poruka u redu)
    \item \textbf{Redis caching} -- smanjenje opterećenja baze podataka
    \item \textbf{CDN (Cloudflare)} -- keširanje statičkog sadržaja bliže korisnicima
    \item \textbf{Connection pooling (PgBouncer)} -- optimizacija konekcija na bazu
\end{itemize}

\textbf{Ciljane metrike:}
\begin{itemize}
    \item Latencija standardnih UI operacija: $<2$ sekunde
    \item Latencija kompleksnih operacija (analitika, izvještaji): $<5$ sekundi
    \item Minimalno $\geq$50 simultanih korisnika bez degradacije performansi
    \item Podrška za $\geq$1.500 krava u sistemu
    \item IoT senzorski podaci: očitavanje svakih 5 minuta, gubitak paketa $<1\%$
\end{itemize}

\subsubsection{Sigurnost}

\textbf{Zaštita mreže:}
\begin{itemize}
    \item \textbf{Tailscale VPN} -- WireGuard-bazirana mesh mreža za sigurnu komunikaciju između čvorova klastera i remote pristup
    \item \textbf{Cloudflare} -- DDoS protekcija, WAF (Web Application Firewall), bot protection
    \item \textbf{TLS 1.2+} -- enkripcija saobraćaja u tranzitu (HTTPS), certifikati putem Let's Encrypt
    \item \textbf{HSTS (HTTP Strict Transport Security)} -- forsiranje HTTPS konekcija
    \item \textbf{mTLS} -- međusobna autentifikacija servisa unutar klastera
\end{itemize}

\textbf{Kontrola pristupa:}
\begin{itemize}
    \item \textbf{Autentifikacija:} JWT tokeni, obavezan 2FA (TOTP) za administratore, OAuth2 za integracije
    \item \textbf{Password politika:} minimalno 12 znakova, kompleksnost, hashiranje sa Argon2 ili bcrypt ($\geq$10 rounds)
    \item \textbf{Session timeout:} 30 minuta neaktivnosti
    \item \textbf{Autorizacija:} Role-Based Access Control (RBAC) -- Administrator, Menadžer, Veterinar, Osoblje
    \item \textbf{API ključevi} -- za M2M (machine-to-machine) komunikaciju sa IoT uređajima, sa rate-limiting
\end{itemize}

\textbf{Zaštita podataka:}
\begin{itemize}
    \item \textbf{Self-hosting} -- svi podaci ostaju pod kontrolom farme, bez dijeljenja sa trećim stranama
    \item \textbf{Enkripcija at-rest} -- AES-256 za podatke u mirovanju, LUKS enkripcija diskova
    \item \textbf{Video snimci} -- retencija 30 dana, nakon čega se automatski brišu
    \item \textbf{Audit trail} -- čuvanje logova pristupa i promjena minimalno 24 mjeseca
    \item \textbf{GDPR usklađenost} -- privacy by design, pravo na brisanje, transparentnost obrade
    \item \textbf{Secrets management} -- Kubernetes Secrets + Sealed Secrets za GitOps
\end{itemize}

\textbf{Backup i oporavak:}
\begin{itemize}
    \item \textbf{Strategija:} Dnevni full backup + inkrementalni backup
    \item \textbf{RPO (Recovery Point Objective):} $\leq$15 minuta
    \item \textbf{RTO (Recovery Time Objective):} $\leq$4 sata
    \item \textbf{Verifikacija:} Automatska provjera integriteta backup-a
    \item \textbf{Off-site replikacija:} Backup na odvojenu lokaciju (Hetzner Storage Box)
    \item \textbf{DR testiranje:} Kvartalno testiranje procedure oporavka
\end{itemize}

\textbf{ALE proračun za procjenu rizika (primjer -- gubitak podataka):}
\begin{align*}
    AV &= 500.000 \text{ KM (vrijednost godišnjih podataka o proizvodnji)} \\
    EF &= 30\% \text{ (procijenjeni gubitak bez backupa)} \\
    SLE &= AV \times EF = 150.000 \text{ KM} \\
    ARO &= 0.1 \text{ (jednom u 10 godina)} \\
    ALE &= SLE \times ARO = 15.000 \text{ KM/godišnje}
\end{align*}

Implementacija backup rješenja (trošak $\approx$ 3.000 KM/god) smanjuje EF na 5\%, čime:
\begin{align*}
    ALE_{nova} &= 500.000 \times 0.05 \times 0.1 = 2.500 \text{ KM/godišnje} \\
    \text{Vrijednost mjere} &= 15.000 - 3.000 - 2.500 = 9.500 \text{ KM/godišnje uštede}
\end{align*}

\subsubsection{Kulturalni i regulatorni zahtjevi}

\begin{itemize}
    \item \textbf{Jezik:} Bosanski/Hrvatski/Srpski (bs/hr/sr) i engleski (en)
    \item \textbf{Lokalizacija:} Lokalizirani datumi, vremenske zone, mjerne jedinice i valute
    \item \textbf{Valuta:} KM (Konvertibilna marka)
    \item \textbf{GDPR usklađenost:} Zaštita ličnih podataka zaposlenih
    \item \textbf{Veterinarski propisi BiH:} Usklađenost sa evidencijama koje zahtijeva veterinarska inspekcija
\end{itemize}

\subsubsection{Pristupačnost i kompatibilnost}

\begin{itemize}
    \item \textbf{WCAG 2.1 AA:} Web klijent ispunjava standarde pristupačnosti (kontrast, navigacija tastaturom, alt opisi)
    \item \textbf{Podržani preglednici:} Zadnje 2 verzije Chrome, Edge, Firefox i Safari
    \item \textbf{Responzivni dizajn:} Podržane rezolucije 360px -- 1920px
    \item \textbf{Mobile-first pristup:} Optimizacija za korištenje na mobilnim uređajima u štali
\end{itemize}

\subsubsection{Offline režim}

\begin{itemize}
    \item \textbf{Lokalno keširanje:} Klijentska aplikacija kešira zapise za rad bez mreže
    \item \textbf{Osnovne funkcije offline:} Unos podataka, pregled zadataka, evidencija muže
    \item \textbf{Automatska sinhronizacija:} Po povratku mreže, podaci se sinhronizuju u $\leq$15 minuta
    \item \textbf{Rješavanje konflikata:} Pregled i manualno rješavanje konfliktnih promjena
    \item \textbf{Edge Gateway:} Lokalni server na farmi kao buffer za IoT podatke pri gubitku internet konekcije
\end{itemize}

\subsection{Dostupnost sistema}

\begin{itemize}
    \item \textbf{Ciljani uptime:} $\geq$99\% mjesečno
    \item \textbf{High Availability:} K3s klaster sa 3 control-plane čvora eliminiše single point of failure
    \item \textbf{Planirani prekidi:} Najava minimalno 48 sati unaprijed
    \item \textbf{Zero-downtime deployment:} Rolling updates bez prekida servisa
\end{itemize}

\subsection{Monitoring i observabilnost}

\begin{itemize}
    \item \textbf{Centralizirani logovi:} Loki za agregaciju i pretragu logova svih servisa
    \item \textbf{Metrike:} Prometheus za prikupljanje metrika (CPU, memorija, latencija, greške)
    \item \textbf{Vizualizacija:} Grafana dashboardi za real-time praćenje sistema
    \item \textbf{Distribuirano praćenje:} OpenTelemetry sa korelacijom trace-ID kroz sve servise
    \item \textbf{Error tracking:} Sentry za automatsko prijavljivanje i alerting aplikacijskih grešaka
    \item \textbf{Uptime monitoring:} Uptime Kuma za praćenje dostupnosti servisa
    \item \textbf{Alerting:} Automatske notifikacije pri prekoračenju pragova (CPU $>$80\%, latencija $>$2s, error rate $>$1\%)
\end{itemize}

\newpage
\section{SPECIFIKACIJA HARDVERA}

Hardverska infrastruktura sistema „OIS -- Smart Cow Farm" podijeljena je na \textbf{cloud infrastrukturu} (Hetzner Cloud) i \textbf{lokalnu infrastrukturu} na farmi.

\subsection{Cloud infrastruktura (Hetzner Cloud)}

\subsubsection{Kubernetes Control Plane čvorovi (3x)}

\begin{center}
\begin{tabular}{|l|l|}
\hline
\rowcolor{lightgray}
\textbf{Komponenta} & \textbf{Specifikacija} \\
\hline
Server tip & Hetzner CPX21 \\
\hline
Procesor & 3 vCPU (AMD EPYC) \\
\hline
RAM & 4 GB DDR4 \\
\hline
Storage & 80 GB NVMe SSD \\
\hline
Mreža & 1 Gbps, 20 TB saobraćaja/mjesec \\
\hline
Lokacija & Falkenstein, Njemačka (EU) \\
\hline
\end{tabular}
\end{center}

\subsubsection{Kubernetes Worker čvorovi -- General Purpose (2-5x, autoscaling)}

\begin{center}
\begin{tabular}{|l|l|}
\hline
\rowcolor{lightgray}
\textbf{Komponenta} & \textbf{Specifikacija} \\
\hline
Server tip & Hetzner CPX41 \\
\hline
Procesor & 8 vCPU (AMD EPYC) \\
\hline
RAM & 16 GB DDR4 \\
\hline
Storage & 240 GB NVMe SSD \\
\hline
Mreža & 1 Gbps \\
\hline
Namjena & API servisi, web aplikacija, baze podataka \\
\hline
\end{tabular}
\end{center}

\subsubsection{Kubernetes Worker čvorovi -- GPU (1-2x, za AI workloade)}

\begin{center}
\begin{tabular}{|l|l|}
\hline
\rowcolor{lightgray}
\textbf{Komponenta} & \textbf{Specifikacija} \\
\hline
Server tip & Hetzner GX11 (GPU instance) \\
\hline
Procesor & 8 vCPU (AMD EPYC) \\
\hline
RAM & 32 GB DDR4 \\
\hline
GPU & NVIDIA A10 (24 GB VRAM) ili RTX 4000 \\
\hline
Storage & 320 GB NVMe SSD \\
\hline
Namjena & AI inference (prepoznavanje krava, analiza ponašanja) \\
\hline
\end{tabular}
\end{center}

\subsubsection{Dedicated Storage Server}

\begin{center}
\begin{tabular}{|l|l|}
\hline
\rowcolor{lightgray}
\textbf{Komponenta} & \textbf{Specifikacija} \\
\hline
Server tip & Hetzner Storage Box BX20 \\
\hline
Kapacitet & 2 TB \\
\hline
Protokoli & SFTP, SCP, rsync, Samba/CIFS \\
\hline
Backup & Automatski snapshoti \\
\hline
Namjena & Backup, arhiva video snimaka, off-site replikacija \\
\hline
\end{tabular}
\end{center}

\subsection{Lokalna infrastruktura na farmi}

\subsubsection{Edge Gateway Server}

\begin{center}
\begin{tabular}{|l|l|}
\hline
\rowcolor{lightgray}
\textbf{Komponenta} & \textbf{Specifikacija} \\
\hline
Uređaj & Intel NUC 12 Pro ili Raspberry Pi 5 (8GB) \\
\hline
Procesor & Intel i5-1240P / ARM Cortex-A76 \\
\hline
RAM & 16 GB / 8 GB \\
\hline
Storage & 512 GB NVMe SSD \\
\hline
Mreža & Gigabit Ethernet, WiFi 6 \\
\hline
Namjena & MQTT broker, lokalno keširanje, edge AI inference \\
\hline
\end{tabular}
\end{center}

\subsubsection{IP kamere za video nadzor (po potrebi)}

\begin{center}
\begin{tabular}{|l|l|}
\hline
\rowcolor{lightgray}
\textbf{Komponenta} & \textbf{Specifikacija} \\
\hline
Model & Hikvision DS-2CD2T47G2 ili slično \\
\hline
Rezolucija & 4MP (2688x1520) \\
\hline
Frame rate & 25 fps \\
\hline
Protokoli & RTSP, ONVIF \\
\hline
IR domet & 80m (za noćni rad u štali) \\
\hline
Zaštita & IP67, vandal-proof \\
\hline
Količina & 4-8 kamera (ovisno o veličini štale) \\
\hline
\end{tabular}
\end{center}

\subsubsection{IoT senzori}

\begin{center}
\begin{tabular}{|l|l|l|}
\hline
\rowcolor{lightgray}
\textbf{Tip senzora} & \textbf{Model} & \textbf{Količina} \\
\hline
Temperatura/Vlaga & DHT22 / SHT31 & 4-8 po štali \\
\hline
Kvalitet zraka (CO2, NH3) & MQ-135 / CCS811 & 2-4 po štali \\
\hline
Mjerač protoka mlijeka & Flowmeter sa RS485 & 1 po muznoj jedinici \\
\hline
Gateway za senzore & ESP32 / Raspberry Pi Pico W & po potrebi \\
\hline
\end{tabular}
\end{center}

\subsubsection{Mrežna oprema na farmi}

\begin{center}
\begin{tabular}{|l|l|}
\hline
\rowcolor{lightgray}
\textbf{Komponenta} & \textbf{Specifikacija} \\
\hline
Router & Mikrotik RB4011 ili Ubiquiti EdgeRouter \\
\hline
PoE Switch & Ubiquiti USW-24-POE (za napajanje kamera) \\
\hline
WiFi Access Point & Ubiquiti U6-LR (za pokrivenost štale) \\
\hline
UPS & APC Smart-UPS 1500VA (za edge server i mrežu) \\
\hline
Internet konekcija & Min. 50 Mbps upload (za video streaming) \\
\hline
\end{tabular}
\end{center}

\newpage
\section{SPECIFIKACIJA SOFTVERA}

\subsection{Operativni sistemi}

\begin{center}
\begin{tabular}{|l|l|l|}
\hline
\rowcolor{lightgray}
\textbf{Komponenta} & \textbf{OS} & \textbf{Verzija} \\
\hline
Cloud serveri & Ubuntu Server LTS & 24.04 \\
\hline
Edge Gateway & Ubuntu Server / Raspberry Pi OS & 24.04 / Bookworm \\
\hline
Kubernetes & K3s & v1.29+ \\
\hline
Kontejneri & Alpine Linux / Debian Slim & latest \\
\hline
\end{tabular}
\end{center}

\subsection{Kontejnerizacija i orkestracija}

\begin{center}
\begin{tabular}{|l|l|l|}
\hline
\rowcolor{lightgray}
\textbf{Softver} & \textbf{Verzija} & \textbf{Namjena} \\
\hline
Docker & 25.x & Kontejnerizacija aplikacija \\
\hline
K3s & 1.29+ & Lightweight Kubernetes distribucija \\
\hline
Helm & 3.x & Package manager za Kubernetes \\
\hline
Traefik & 3.x & Ingress controller, reverse proxy \\
\hline
Cert-Manager & 1.x & Automatsko izdavanje TLS certifikata \\
\hline
\end{tabular}
\end{center}

\subsection{Baze podataka i skladištenje}

\begin{center}
\begin{tabular}{|l|l|l|}
\hline
\rowcolor{lightgray}
\textbf{Softver} & \textbf{Verzija} & \textbf{Namjena} \\
\hline
PostgreSQL & 16.x & Primarna relaciona baza podataka \\
\hline
TimescaleDB & 2.x & Ekstenzija za vremenske serije (senzori) \\
\hline
Redis & 7.x & Caching, session storage, pub/sub \\
\hline
MinIO & latest & S3-kompatibilno objektno skladište \\
\hline
PgBouncer & 1.x & Connection pooling za PostgreSQL \\
\hline
\end{tabular}
\end{center}

\subsection{Backend tehnologije}

\begin{center}
\begin{tabular}{|l|l|l|}
\hline
\rowcolor{lightgray}
\textbf{Tehnologija} & \textbf{Verzija} & \textbf{Namjena} \\
\hline
Python & 3.12+ & AI servisi, data processing \\
\hline
FastAPI & 0.110+ & REST API framework (Python) \\
\hline
Node.js & 20 LTS & Real-time servisi, Notification service \\
\hline
Go & 1.22+ & High-performance servisi (opciono) \\
\hline
\end{tabular}
\end{center}

\subsection{Frontend tehnologije}

\begin{center}
\begin{tabular}{|l|l|l|}
\hline
\rowcolor{lightgray}
\textbf{Tehnologija} & \textbf{Verzija} & \textbf{Namjena} \\
\hline
React / Vue.js & 18.x / 3.x & SPA web aplikacija \\
\hline
TypeScript & 5.x & Type-safe JavaScript \\
\hline
TailwindCSS & 3.x & Utility-first CSS framework \\
\hline
Chart.js / D3.js & latest & Vizualizacija podataka \\
\hline
\end{tabular}
\end{center}

\subsection{AI/ML stack}

\begin{center}
\begin{tabular}{|l|l|l|}
\hline
\rowcolor{lightgray}
\textbf{Softver} & \textbf{Verzija} & \textbf{Namjena} \\
\hline
PyTorch & 2.x & Deep learning framework \\
\hline
ONNX Runtime & 1.x & Optimizirano izvršavanje modela \\
\hline
NVIDIA Triton & 24.x & Model serving infrastruktura \\
\hline
OpenCV & 4.x & Obrada video streama \\
\hline
Ultralytics YOLO & v8 & Object detection (prepoznavanje krava) \\
\hline
CUDA Toolkit & 12.x & GPU akceleracija \\
\hline
\end{tabular}
\end{center}

\subsection{IoT i messaging}

\begin{center}
\begin{tabular}{|l|l|l|}
\hline
\rowcolor{lightgray}
\textbf{Softver} & \textbf{Verzija} & \textbf{Namjena} \\
\hline
EMQX & 5.x & MQTT broker za IoT senzore \\
\hline
NATS & 2.x & Lightweight messaging između servisa \\
\hline
\end{tabular}
\end{center}

\subsection{Sigurnost i mreža}

\begin{center}
\begin{tabular}{|l|l|l|}
\hline
\rowcolor{lightgray}
\textbf{Softver} & \textbf{Verzija} & \textbf{Namjena} \\
\hline
Tailscale & latest & WireGuard VPN mesh mreža \\
\hline
Cloudflare & -- & DDoS zaštita, WAF, CDN \\
\hline
Let's Encrypt & -- & Besplatni TLS certifikati \\
\hline
Authelia / Keycloak & latest & Identity provider, SSO (opciono) \\
\hline
\end{tabular}
\end{center}

\subsection{Monitoring i observabilnost}

\begin{center}
\begin{tabular}{|l|l|l|}
\hline
\rowcolor{lightgray}
\textbf{Softver} & \textbf{Verzija} & \textbf{Namjena} \\
\hline
Prometheus & 2.x & Prikupljanje metrika \\
\hline
Grafana & 10.x & Vizualizacija i dashboardi \\
\hline
Loki & 2.x & Agregacija logova \\
\hline
OpenTelemetry & latest & Distribuirano praćenje (tracing) \\
\hline
Sentry & latest & Error tracking i alerting \\
\hline
Uptime Kuma & latest & Status page i uptime monitoring \\
\hline
\end{tabular}
\end{center}

\subsection{DevOps i CI/CD}

\begin{center}
\begin{tabular}{|l|l|l|}
\hline
\rowcolor{lightgray}
\textbf{Softver} & \textbf{Verzija} & \textbf{Namjena} \\
\hline
Git & 2.x & Version control \\
\hline
GitHub Actions / GitLab CI & -- & CI/CD pipeline \\
\hline
ArgoCD & 2.x & GitOps deployment na Kubernetes \\
\hline
Terraform & 1.x & Infrastructure as Code \\
\hline
Ansible & 2.x & Konfiguracija servera \\
\hline
SonarQube & latest & Statička analiza koda \\
\hline
pytest / Jest & latest & Unit i integration testovi \\
\hline
\end{tabular}
\end{center}

\subsection{Backup i disaster recovery}

\begin{center}
\begin{tabular}{|l|l|l|}
\hline
\rowcolor{lightgray}
\textbf{Softver} & \textbf{Verzija} & \textbf{Namjena} \\
\hline
Velero & 1.x & Kubernetes backup \\
\hline
pgBackRest & 2.x & PostgreSQL backup i PITR \\
\hline
Restic & 0.16+ & Enkriptovani backup na remote storage \\
\hline
\end{tabular}
\end{center}

\end{document}